\documentclass[lineno]{JFM-FLM_Au}

\usepackage{xcolor}
\usepackage{tikz}
\usetikzlibrary{calc,arrows.meta}
\usepackage{subcaption}

\newcommand{\tbd}[1]{\textcolor{red}{\ensuremath{\langle\mbox{\small #1}\rangle}}}
\definecolor{jfmblue}{RGB}{0,56,168}
\def\xshift{1.5}

\hypersetup{
  colorlinks=true,
  citecolor=blue,
  linkcolor=jfmblue,
  urlcolor=jfmblue
}

\newtheorem{lemma}{Lemma}
\newtheorem{corollary}{Corollary}
\newtheorem{proposition}{Proposition}

\newenvironment{proof}
  {\par\noindent\textit{Proof.}\ }
  {\hfill$\Box$\par}

\newcounter{assumption}
\renewcommand{\theassumption}{A\arabic{assumption}}

\newenvironment{assumption}[1]{%
  \refstepcounter{assumption}%
  \par\medskip
  \noindent\textbf{Assumption~\theassumption\ (#1).}\ %
}{%
  \par\medskip
}  

\lefttitle{Sharma et al.}
\righttitle{Journal of Fluid Mechanics}

\title{Strain-coupled one-dimensional turbulence for rapid distortion}

\author{Sparsh Sharma\aff{1}}

\affiliation{\aff{1}German Aerospace Center (DLR), 38108 Braunschweig, Germany}

\corresau{Sparsh Sharma, sparsh.sharma@dlr.de}

\begin{document}
\maketitle

\begin{abstract}
xxxx
\end{abstract}

\begin{keywords}
Authors should not enter keywords on the manuscript, as these must be chosen by the author during the online submission process and will then be added during the typesetting process (see \href{https://www.cambridge.org/core/services/aop-file-manager/file/61436b61ff7f3cfab749ce3a/JFM-Keywords-Sept-2021.pdf.}{Keyword PDF} for the full list).  Other classifications will be added at the same time.
\end{keywords}

\section{Introduction}\label{sec:intro}

Turbulence--airfoil interaction noise, radiated when a lifting surface
encounters oncoming turbulence, is a dominant broadband source in axial
fans, contra-rotating rotors, wind-turbine and propeller blades, and
airframe components operating in disturbed inflow. When the integral
length scale of the incident turbulence is large compared with the
leading-edge radius, the radiated field is produced predominantly at the
leading edge, where the unsteady upwash induces a fluctuating surface
loading that scatters into sound as a non-compact dipole. Predicting this
leading-edge noise both reliably and at a cost compatible with design
iteration remains an open problem, and is the practical motivation for the
present work.

The prevailing predictive framework is the analytical theory of
\citet{Amiet1975}, which combines the linearised gust-response solution of
\citet{Sears1941} with the surface-source formulation of \citet{Curle1955}
and can be embedded within the general moving-surface analogy of
\citet{FfowcsWilliams1969}. In Amiet's theory the far-field acoustic
power spectral density is obtained as a linear functional of the
wavenumber--frequency spectrum of the upwash (airfoil-normal) velocity
component of the incoming turbulence, convolved with an airfoil response
function and an edge-scattering kernel and weighted by the spanwise
correlation of the loading. The framework has been extended to finite
chord, finite span and realistic geometry and validated against
grid-turbulence experiments \citep{PatersonAmiet1977,RogerMoreau2010},
but the structure of the prediction is unchanged: \emph{the upwash
spectrum that enters the theory is prescribed, not computed}. In practice
it is taken from an idealised homogeneous, isotropic model spectrum---most
commonly the von~K\'arm\'an form \citep{vonKarman1948}---fixed by the
upstream turbulence intensity and integral length scale and frozen in
space through Taylor's hypothesis. Because the radiated spectrum is linear
in this input, the fidelity of any leading-edge-noise prediction is
limited, first and foremost, by the fidelity of the prescribed upwash
spectrum.

The assumption that the turbulence reaching the leading edge is identical
to the undisturbed upstream turbulence is, however, physically
questionable. As the flow approaches the stagnation region the mean
velocity gradient distorts the incident turbulence over a time that, for
the energy-containing scales, can be comparable to or shorter than the
eddy-turnover time. This is the rapid-distortion regime first analysed by
\citet{BatchelorProudman1954} and developed for flow around obstacles by
\citet{Hunt1973} and \citet{Goldstein1978}: the turbulence becomes
anisotropic and its spectral content is reorganised, while the body
imposes a kinematic ``blocking'' that suppresses the wall-normal velocity
component close to the surface even as the straining amplifies it farther
upstream \citep{HuntGraham1978,GoldsteinAtassi1976}. The upwash spectrum
\emph{seen by the leading edge} is therefore distorted, scale dependent
and tensorially redistributed relative to its upstream counterpart. This
pre-impact distortion stage lies precisely between the upstream
measurement plane and the input required by Amiet's theory, and it is the
stage that the standard prediction chain omits.

Several routes to recovering this distortion exist, each with a
characteristic limitation. Scale-resolving simulation---large-eddy
simulation, lattice-Boltzmann methods or direct computational
aeroacoustics---represents the distortion and the radiation directly, but
at a cost that precludes broad parametric study across turbulence
intensity, length-scale ratio and Reynolds number. Linear rapid-distortion
theory (RDT) provides an analytical description of the mean-strain
distortion and, in its inhomogeneous form, of the surface blocking
\citep{HuntCarruthers1990,GoldsteinAfsarLeib2013}; it has been used to
construct RDT-modified turbulence spectra that are fed to Amiet's solution
in place of the isotropic input \citep{deSantana2016}. RDT is, however,
linear, and in its tractable homogeneous form it neglects the nonlinear
inter-scale transfer and the slow return-to-isotropy that act over the
finite approach time. Synthetic-turbulence methods---synthetic-eddy and
random-particle techniques---efficiently reproduce a \emph{prescribed}
spectrum and coherence and are now routine for generating aeroacoustic
inflow \citep{KimHaeri2015}, but by construction they impose the spectrum
rather than evolving it, and so cannot themselves predict how the spectrum
is modified before impact. What is missing is a reduced-order model that
\emph{evolves} the incoming turbulence through the distortion region with
full scale resolution and at negligible cost, delivering the distorted
upwash spectrum as an output rather than an assumption.

The one-dimensional turbulence (ODT) model of \citet{Kerstein1999} is a
natural candidate for this role. ODT represents the evolution of a
turbulent velocity profile along a single notional line of sight,
resolving the full range of length scales in one spatial dimension while
modelling turbulent advection as a stochastic sequence of scale-local
mapping (``eddy'') events built on the triplet map introduced in the
linear-eddy model \citep{Kerstein1991}. Its vector formulation
\citep{Kerstein2001} carries all three velocity components and
redistributes energy among them through a kernel mechanism that emulates
pressure scrambling \citep{WunschKerstein2001,Kerstein2001}, reproducing
inter-component transfer and return to isotropy; the model has been generalised to variable density
\citep{AshurstKerstein2005}, to efficient adaptive-mesh and spatially
developing implementations \citep{Lignell2013,StephensLignell2021}, to
cylindrical and spherical geometries \citep{Lignell2018}, and to a
three-dimensional large-eddy framework \citep{SchmidtKerstein2010}. ODT
reproduces inertial-range scaling and component statistics of free-shear
and wall-bounded flows at a small fraction of the cost of
three-dimensional simulation. Of particular relevance here, ODT velocity
fields have recently been coupled to acoustic analogies to estimate the
far-field sound of low-Mach-number turbulent jets
\citep{SharmaKleinSchmidt2022,MedinaMendez2023}, demonstrating that ODT
statistics---including the component spectra and fluctuating Reynolds
stresses---can be supplied to an external acoustic theory. These
developments establish ODT as an accurate and inexpensive generator of
scale-resolved turbulence statistics, and in particular of the
wall-normal velocity spectrum required by Amiet's theory.

Standard ODT cannot, however, represent the pre-impact distortion itself.
The model contains no mechanism for the action of an imposed mean strain
on the fluctuating field: its pressure treatment, embodied in the
scrambling kernel, models only the slow, turbulence-driven return to
isotropy and not the rapid, mean-strain-driven redistribution that governs
the stagnation region. The apparent remedy---importing the linear RDT
amplitude equation, in which the pressure response is the solenoidal
projection $P_{ij}(\boldsymbol{k}) = \delta_{ij} - k_i k_j/k^2$ acting in
three-dimensional wavenumber space---is fundamentally incompatible with
the one-dimensional model, because the projection requires the full
wavevector $\boldsymbol{k}=(k_1,k_2,k_3)$, whereas a single ODT line
resolves only the wavenumber along its own coordinate. Reconciling the
inherently three-dimensional and non-local pressure physics of rapid
distortion with the one-dimensional and local structure of ODT is the
central modelling problem addressed in this paper.

We present a strain-coupled formulation of one-dimensional turbulence that
incorporates rapid distortion by an imposed mean velocity gradient while
preserving the structural foundations of the model. The formulation rests
on a decomposition of the pressure--strain interaction into its rapid and
slow parts \citep{Pope2000}: the slow, turbulence-driven part is retained
as the existing ODT scrambling kernel, while the rapid, mean-strain-driven
part is introduced as a new component-redistribution kernel. Rather than
evaluating the three-dimensional projection, the rapid kernel is
constructed so that the component-energy evolution of the line reproduces
the exact homogeneous-RDT solution for an imposed mean strain, rendering
the formulation provably RDT-consistent in the homogeneous limit while
remaining expressed entirely in quantities available on the line. The
mean-strain production enters as a continuous forcing of the velocity
components, which we show preserves the measure-preservation, continuity
and conservation properties of the triplet-map mechanism; the kinematic
compression of wavenumbers under strain is captured by a consistent
deformation of the ODT domain, connecting the formulation to its
adaptive-mesh implementation. The result is a reduced-order, scale-resolved
model of the pre-impact distortion stage whose output---the distorted,
anisotropic upwash spectrum---is intended to replace the prescribed
isotropic input of Amiet-type prediction, with the airfoil response, edge
scattering and spanwise coherence retained as the classical, downstream
acoustic stage.

The remainder of the paper is organised as follows.
Section~\ref{sec:odt} reviews the elements of standard ODT required for the
present development and the rapid/slow structure of the pressure--strain
interaction. Section~\ref{sec:formulation} derives the strain-coupled
formulation, establishes the preservation of the model's structural
invariants under mean-strain forcing, and constructs the RDT-consistent
rapid kernel. Section~\ref{sec:validation} validates the formulation
against analytical homogeneous-RDT solutions for the canonical strains
relevant to the stagnation region---axisymmetric contraction and plane
strain---and against reference scale-resolving data. Section~\ref{sec:upwash}
demonstrates the resulting distorted upwash spectrum and its implications
for the Amiet input. Section~\ref{sec:conclusions} concludes and outlines
the remaining steps---surface blocking and spanwise coherence---towards a
complete leading-edge-noise framework.

% =====================================================================
%  Section 2  --  Background
%  (i) standard one-dimensional turbulence; (ii) rapid/slow structure
%      of the pressure--strain interaction and linear RDT.
%  Target: Journal of Fluid Mechanics (jfm.cls); citations in references.bib
% =====================================================================

\section{Background}\label{sec:odt}

This section collects the elements on which the strain-coupled formulation
of \S\,\ref{sec:formulation} is built. Section~\ref{sec:odt-standard}
summarises standard one-dimensional turbulence (ODT) in the notation used
throughout; \S\,\ref{sec:rapid-slow} recalls the rapid/slow decomposition
of the pressure--strain interaction and the linear rapid-distortion
equations that the new formulation is required to reproduce. Nothing in
this section is original; readers familiar with
\citet{Kerstein1999,Kerstein2001} and with second-moment closure
\citep{Pope2000} may proceed to \S\,\ref{sec:formulation} after noting the
notation and the correspondence stated in \S\,\ref{sec:rapid-slow}.

\subsection{Standard one-dimensional turbulence}\label{sec:odt-standard}

ODT evolves a three-component velocity field $u_i(y,t)$, $i=1,2,3$, and any
number of scalar fields $\theta(y,t)$, defined on a one-dimensional domain
parameterised by the spatial coordinate $y$, which is identified with the
coordinate direction $x_2$ \citep{Kerstein2001}. The fields represent
individual realisations: ensemble statistics are accumulated over many
independent simulations rather than carried by the fields themselves. The
fields evolve by two interleaved mechanisms---continuous molecular
transport and a stochastic sequence of instantaneous \emph{eddy events}.

\subsubsection{Molecular transport}

Between events the fields obey ordinary one-dimensional diffusion,
\begin{equation}
  \left(\partial_t-\nu\,\partial_y^2\right)u_i(y,t)=0,
  \qquad
  \left(\partial_t-\kappa\,\partial_y^2\right)\theta(y,t)=0,
  \label{eq:diffusion}
\end{equation}
with kinematic viscosity $\nu$ and scalar diffusivity $\kappa$
(Schmidt number $\mathit{Sc}=\nu/\kappa$). In a constant-property,
constant-density medium this is the only continuous evolution; the
variable-density generalisation is given by \citet{AshurstKerstein2005}.

\subsubsection{Eddy events: the triplet map and kernel}

An eddy event models the effect of a single turbulent overturn. Because a
continuum rotational motion cannot be represented on a line, the event is
implemented as an instantaneous, measure-preserving rearrangement of the
profiles together with an additive velocity change,
\begin{equation}
  u_i(y)\;\longrightarrow\;u_i\!\big(f(y)\big)+c_i\,K(y),
  \qquad
  \theta(y)\;\longrightarrow\;\theta\!\big(f(y)\big).
  \label{eq:event}
\end{equation}
The map $f$ is the \emph{triplet map} \citep{Kerstein1991,Kerstein1999},
defined on an interval $[y_0,y_0+l]$ of size $l$ by
\begin{equation}
  f(y)=y_0+
  \begin{cases}
    3(y-y_0), & y_0\le y\le y_0+\tfrac{1}{3}l,\\[2pt]
    2l-3(y-y_0), & y_0+\tfrac{1}{3}l\le y\le y_0+\tfrac{2}{3}l,\\[2pt]
    3(y-y_0)-2l, & y_0+\tfrac{2}{3}l\le y\le y_0+l,\\[2pt]
    y-y_0, & \text{otherwise,}
  \end{cases}
  \label{eq:triplet}
\end{equation}
written here through its inverse, so that fluid at $f(y)$ is placed at $y$.
The map compresses the interval to one-third of its length, inserts three
copies, and reverses the central copy; it leaves the field outside
$[y_0,y_0+l]$ unchanged. The triplet map satisfies the three structural
requirements that underpin the model: it is \emph{measure preserving} (the
non-local analogue of a divergence-free rearrangement), \emph{continuous}
(it introduces no discontinuities into the fields), and \emph{scale local}
(it changes property gradients by at most a factor of order unity)
\citep{Kerstein2001}. These three properties are exactly the invariants
whose preservation under an imposed mean strain is established in
\S\,\ref{sec:formulation}.

The additive term in \eqref{eq:event} is built from the \emph{kernel}
\begin{equation}
  K(y)\equiv y-f(y),
  \label{eq:kernel}
\end{equation}
which is non-zero only within the eddy interval and satisfies the two
identities used below,
\begin{equation}
  \int K(y)\,\mathrm{d}y = 0,
  \qquad
  \int K^2(y)\,\mathrm{d}y = \tfrac{4}{27}\,l^{3}.
  \label{eq:kernel-identities}
\end{equation}
The first guarantees that the additive change conserves the $y$-integrated
momentum of each velocity component at constant density; the second fixes
the energetics of the redistribution. The amplitudes $c_i$ are determined
event-by-event by the pressure-scrambling rule described next.

\subsubsection{Pressure scrambling and component-energy redistribution}

The additive kernel term in \eqref{eq:event} is the model surrogate for the
pressure-mediated redistribution of energy among velocity components, an
effect absent from the original single-component formulation
\citep{Kerstein1999} and introduced through the kernel by
\citet{WunschKerstein2001,Kerstein2001}. Using measure preservation and
\eqref{eq:kernel-identities}, the change in the kinetic energy of component
$i$ produced by an event is
\begin{equation}
  \Delta E_i = \rho_0\, l^2\, c_i\!\left(u_{i,K}+\tfrac{2}{27}\,l\,c_i\right),
  \qquad
  u_{i,K}\equiv\frac{1}{l^2}\int u_i\!\big(f(y)\big)\,K(y)\,\mathrm{d}y,
  \label{eq:dEi}
\end{equation}
where $\rho_0$ is the (constant) density per unit length, which does not
affect the constant-density applications considered here. The amplitudes
are fixed by requiring that the event redistribute energy among components
while conserving the total, $\sum_i\Delta E_i=0$, and that the rule be
invariant under permutation of the component indices. These requirements
give
\begin{equation}
  \Delta E_i = \alpha\sum_j T_{ij}\,Q_j,
  \qquad
  \mathsf{T}=\frac{1}{2}
  \begin{pmatrix}
    -2 & 1 & 1\\ 1 & -2 & 1\\ 1 & 1 & -2
  \end{pmatrix},
  \qquad
  Q_i\equiv\tfrac{27}{8}\,\rho_0\, l\, u_{i,K}^2,
  \label{eq:redistribution}
\end{equation}
in which $Q_i$ is the \emph{available kinetic energy} of component $i$
within the eddy and $\alpha\in[0,1]$ is the model parameter setting the
fraction of extractable energy transferred towards isotropy
\citep{Kerstein2001}. Combining \eqref{eq:dEi} and
\eqref{eq:redistribution} determines the amplitudes,
\begin{equation}
  c_i = \frac{27}{4l}\left[-u_{i,K}
        +\operatorname{sgn}(u_{i,K})
        \sqrt{u_{i,K}^2+\alpha\sum_j T_{ij}\,u_{j,K}^2}\,\right],
  \label{eq:ci}
\end{equation}
which is real for $0\le\alpha\le1$ and vanishes as $\alpha\to0$. The
transfer matrix $\mathsf{T}$ acts to equalise the component energies; the
construction \eqref{eq:redistribution}--\eqref{eq:ci} is therefore a model
of the \emph{return-to-isotropy} (slow) part of the pressure--strain
interaction, a point made precise in \S\,\ref{sec:rapid-slow}. Standard ODT
contains no counterpart of the rapid, mean-strain-driven part.

\subsubsection{Eddy selection}

The event sequence is a Poisson process whose rate reflects the
instantaneous state of the flow. Each candidate eddy $(y_0,l)$ is assigned
a time scale $\tau(y_0,l;t)$ built from the available energy of the
velocity component aligned with the domain,
\begin{equation}
  \left(\frac{l}{\tau}\right)^{2}
  \sim\; u_{2,K}^2+\alpha\sum_j T_{2j}\,u_{j,K}^2-Z\,\frac{\nu^2}{l^2},
  \label{eq:timescale}
\end{equation}
the final term being a viscous penalty, controlled by the order-unity
parameter $Z$, that suppresses eddies slower than their viscous time. The
corresponding event-rate distribution is
\begin{equation}
  \lambda(y_0,l;t)=\frac{C}{l^2\,\tau(y_0,l;t)},
  \label{eq:rate}
\end{equation}
with rate constant $C$; eddies for which the right-hand side of
\eqref{eq:timescale} is negative are forbidden. Crucially, the
inertial-range cascade and the $E(k)\sim k^{-5/3}$ scaling are
\emph{outcomes} of repeated sampling from \eqref{eq:rate}, not inputs to
it \citep{Kerstein1999}. The three constants $C$, $Z$ and a large-eddy
suppression parameter are the only tunable inputs of the model
\citep{StephensLignell2021}.

\subsubsection{Statistical observables}\label{sec:observables}

ODT yields single- and multi-point moments, probability density functions,
conditional statistics, power spectra and Lagrangian statistics
\citep{Kerstein1999}. For the present work the relevant observable is the
single-component velocity spectrum---in particular the spectrum and
variance of the component normal to the airfoil, which are second moments
of the velocity field and are obtained directly from the realisations. A
distinction emphasised by \citet{Kerstein2001} must, however, be respected:
because the line velocity does not itself advect fluid, an off-diagonal
Reynolds stress such as $\langle u_1' u_2'\rangle$ \emph{cannot} be
interpreted as a momentum flux when formed as a product of the line
velocity components; such advective fluxes are instead evaluated by
monitoring the transport induced by the eddy events. This distinction
constrains how the anisotropy produced by the formulation of
\S\,\ref{sec:formulation} is to be diagnosed and validated
(\S\,\ref{sec:validation}): diagonal component spectra are read from the
fields, whereas shear stresses are read from eddy-induced fluxes.

\subsection{Rapid and slow pressure--strain, and linear rapid
distortion}\label{sec:rapid-slow}

The Reynolds-stress transport equation for a fluctuating field $u_i'$ about
a mean $U_i$ with constant velocity gradient $A_{ij}\equiv\partial
U_i/\partial x_j$ reads, in homogeneous turbulence,
\begin{equation}
  \frac{\mathrm{d}\langle u_i' u_j'\rangle}{\mathrm{d}t}
  = \mathcal{P}_{ij} + \Pi_{ij} - \varepsilon_{ij},
  \label{eq:rstransport}
\end{equation}
with production, pressure--strain and dissipation
\begin{equation}
  \mathcal{P}_{ij}=-\langle u_i' u_k'\rangle A_{jk}
                   -\langle u_j' u_k'\rangle A_{ik},
  \quad
  \Pi_{ij}=\Big\langle \tfrac{p'}{\rho}
            \big(\partial_j u_i'+\partial_i u_j'\big)\Big\rangle,
  \quad
  \varepsilon_{ij}=2\nu\langle \partial_k u_i'\,\partial_k u_j'\rangle .
  \label{eq:Pij}
\end{equation}
Taking the divergence of the fluctuating momentum equation gives the
pressure as the solution of a Poisson equation whose source splits into a
part linear in the mean velocity gradient and a part quadratic in the
fluctuations,
\begin{equation}
  \frac{1}{\rho}\nabla^2 p'
  = \underbrace{-2\,A_{lm}\,\partial_l u_m'}_{\text{rapid}}
    \;-\;
    \underbrace{\partial_l\partial_m\!\left(u_l' u_m'
     -\langle u_l' u_m'\rangle\right)}_{\text{slow}} .
  \label{eq:poisson}
\end{equation}
The pressure--strain correlation inherits this decomposition,
$\Pi_{ij}=\Pi_{ij}^{(r)}+\Pi_{ij}^{(s)}$, into a \emph{rapid} part driven
by the mean strain and a \emph{slow} part driven by the turbulence alone
\citep{Pope2000}. The slow part relaxes the Reynolds stress towards
isotropy; the simplest closure is that of Rotta,
$\Pi_{ij}^{(s)}=-C_1(\varepsilon/k)\big(\langle u_i'u_j'\rangle
-\tfrac{2}{3}k\,\delta_{ij}\big)$, with $k=\tfrac12\langle u_i'u_i'\rangle$.
It is this return-to-isotropy mechanism that the ODT kernel
\eqref{eq:redistribution}--\eqref{eq:ci} reproduces in model form. The
rapid part is linear in $A_{lm}$ and is determined, in homogeneous
turbulence, by the velocity-spectrum tensor
$\Phi_{ij}(\boldsymbol{k})$; the corresponding ``isotropisation of
production'' closure is
$\Pi_{ij}^{(r)}=-C_2\big(\mathcal{P}_{ij}-\tfrac{2}{3}\mathcal{P}\,
\delta_{ij}\big)$, with $\mathcal{P}=\tfrac12\mathcal{P}_{ii}$
\citep{Pope2000}. Standard ODT represents neither $\mathcal{P}_{ij}$ nor
$\Pi_{ij}^{(r)}$.

For a rapidly applied strain---formally, when the total strain accumulates
in a time short compared with the eddy-turnover time---the rapid part can
be computed exactly from linear theory \citep{BatchelorProudman1954,
Hunt1973,HuntCarruthers1990}. Writing the fluctuation as a sum of Fourier
modes that follow the mean flow, each mode of (time-dependent) wavevector
$\boldsymbol{k}(t)$ and amplitude $\hat{u}_i(\boldsymbol{k},t)$ evolves,
in the inviscid rapid limit, according to
\begin{equation}
  \frac{\mathrm{d}k_i}{\mathrm{d}t}=-A_{ji}\,k_j,
  \qquad
  \frac{\mathrm{d}\hat{u}_i}{\mathrm{d}t}
  =\left(\frac{2\,k_i k_m}{k^2}-\delta_{im}\right)A_{mn}\,\hat{u}_n ,
  \label{eq:rdt}
\end{equation}
where $k^2=k_m k_m$. The bracketed operator combines the production
$-A_{im}\hat{u}_m$ with the rapid pressure response, the latter being the
solenoidal projection
$P_{ij}(\boldsymbol{k})=\delta_{ij}-k_ik_j/k^2$ applied to the strained
field. The factor of two in \eqref{eq:rdt} is fixed by the requirement that
the motion remain divergence free: using both relations in
\eqref{eq:rdt},
\begin{equation}
  \frac{\mathrm{d}}{\mathrm{d}t}\big(k_i\hat{u}_i\big)
  = \dot{k}_i\hat{u}_i + k_i\dot{\hat{u}}_i
  = -A_{ji}k_j\hat{u}_i
    + \Big(\tfrac{2k^2 k_m}{k^2}-k_m\Big)A_{mn}\hat{u}_n
  = 0,
  \label{eq:solenoidal}
\end{equation}
so that $k_i\hat{u}_i=0$ is preserved; any other coefficient violates
continuity. Equations~\eqref{eq:rdt} integrate the spectral tensor through
the distortion and thereby determine the rapid pressure--strain
$\Pi_{ij}^{(r)}$ and the component-energy evolution exactly.

Two features of \eqref{eq:rdt} are central to what follows. First, the
pressure response is the operator $P_{ij}(\boldsymbol{k})$, which is
\emph{non-local} (it couples the components through the full wavevector)
and intrinsically \emph{three-dimensional} (it requires all of $k_1,k_2,k_3$).
A single ODT line resolves only the wavenumber along its own coordinate,
so \eqref{eq:rdt} cannot be evaluated on the line as written; reconciling
this incompatibility is the subject of \S\,\ref{sec:formulation}. Second,
the decomposition $\Pi_{ij}=\Pi_{ij}^{(r)}+\Pi_{ij}^{(s)}$ furnishes the
organising principle of the new formulation: the slow part is already
carried by the ODT kernel \eqref{eq:redistribution}--\eqref{eq:ci}, and it
remains only to supply the production $\mathcal{P}_{ij}$ and a rapid
redistribution consistent with \eqref{eq:rdt}, expressed in quantities
available on the line.

% =====================================================================
%  Section 3  --  Strain-coupled one-dimensional turbulence
%  The new formulation: production forcing, RDT-consistent rapid kernel,
%  domain straining. Self-contained derivations.
%
%  Target: Journal of Fluid Mechanics (jfm.cls).
%  Requires in preamble:
%     \usepackage{amsmath,amssymb,amsthm}
%     \newtheorem{proposition}{Proposition}[section]
%  Citations resolve against references.bib.
% =====================================================================

\section{Strain-coupled one-dimensional turbulence}\label{sec:formulation}

The linear rapid-distortion equations~\eqref{eq:rdt} contain two distinct
ingredients: an evolution of the Fourier \emph{amplitudes}
$\hat{u}_i(\boldsymbol{k},t)$, comprising production and the solenoidal
pressure projection, and an evolution of the \emph{wavevectors}
$\boldsymbol{k}(t)$ that carries the kinematic straining of scales. The
amplitude equation is the obstruction identified in
\S\,\ref{sec:rapid-slow}: its pressure term requires the full
three-dimensional wavevector and is therefore not directly representable on
a one-dimensional line. Our strategy is to split the distortion physics
according to this structure and to assign each part to the ODT mechanism
that can carry it faithfully. The production, which is local in physical
space, is imposed as a continuous forcing of the line velocity
(\S\,\ref{sec:strain-forcing}); the rapid pressure--strain, which is not,
is supplied as a continuous, energy-conserving redistribution operator
constructed to reproduce the homogeneous rapid-distortion evolution of the
component energies under an explicit spectral closure
(\S\,\ref{sec:rapid-kernel}); and the wavevector kinematics are recovered by
a consistent straining of the ODT domain
(\S\,\ref{sec:domain-strain}). The slow, turbulence-driven
pressure--strain and the inertial cascade remain the discrete eddy events
of standard ODT, unchanged. Section~\ref{sec:summary} assembles the
formulation and identifies the parameter governing the balance between the
rapid and slow mechanisms.

Throughout, $A_{ij}=\partial U_i/\partial x_j$ denotes the imposed mean
velocity gradient, taken uniform across the line at each instant and
varying along the mean trajectory; the medium is incompressible,
$A_{ii}=0$. The line carries the fluctuation field $u_i(y,t)$ about this
mean. We write $R_{ij}\equiv\overline{u_i u_j}$ for the Reynolds-stress
tensor, where the overline denotes the average over the statistically
homogeneous directions (equivalently, the ensemble of independent
realisations), and $k_t\equiv\tfrac12 R_{ii}$ for the turbulent kinetic
energy. In the homogeneous configurations used to establish and test the
formulation (\S\,\ref{sec:validation}) the spatial line average and the
ensemble average coincide.

\subsection{The strain-coupled line equation and its structural
admissibility}\label{sec:strain-forcing}

In rapid-distortion theory the production term of the fluctuation momentum
equation, $-A_{ij}u_j$, is a pointwise linear combination of the local
velocity components and is therefore exactly representable on the line.
Between eddy events we replace the molecular evolution
\eqref{eq:diffusion} by
\begin{equation}
  \partial_t u_i
  = \nu\,\partial_y^2 u_i \;-\; A_{ij}\,u_j \;+\; B_{ij}\,u_j,
  \label{eq:cont-evolution}
\end{equation}
in which $-A_{ij}u_j$ is the production and $B_{ij}u_j$ is the rapid
pressure--strain operator constructed in \S\,\ref{sec:rapid-kernel}; we
defer $B_{ij}$ and treat the production term first. The eddy events
\eqref{eq:event}--\eqref{eq:ci} are retained without modification. The
evolution is thus advanced by operator splitting between a continuous phase
governed by \eqref{eq:cont-evolution} and a discrete sequence of events
sampled from \eqref{eq:rate}.

A formulation is admissible only if the additional forcing does not
violate the structural properties on which ODT rests. We establish this
directly.

\begin{proposition}[Structural admissibility]\label{prop:admissibility}
The strain forcing in \eqref{eq:cont-evolution} preserves the four
structural properties of standard ODT: \emph{(i)} the measure-preservation,
\emph{(ii)} continuity and \emph{(iii)} scale-locality of the triplet map,
and \emph{(iv)} the kernel identities~\eqref{eq:kernel-identities}.
\end{proposition}

\begin{proof}
Properties (i)--(iv) are properties of the discrete event, that is, of the
map $f$~\eqref{eq:triplet} and the kernel $K$~\eqref{eq:kernel}. Because the
advancement is split into a continuous phase \eqref{eq:cont-evolution} and
the unaltered discrete event, the map and kernel are applied exactly as in
standard ODT; properties (i), (iii) and (iv), which are geometric
statements about $f$ and $K$ on a fixed domain, therefore hold verbatim.
It remains to verify that the continuous phase does not destroy the
continuity of the fields, which underlies (ii). For a mean gradient uniform
along the line, the linear, $y$-independent operator
$\mathcal{A}_{ij}\equiv -A_{ij}+B_{ij}$ acts pointwise, and the
solution of the advective part of \eqref{eq:cont-evolution} over a
sub-step $[t_0,t]$ is
\begin{equation}
  u_i(y,t)=\big[\exp\!\textstyle\int_{t_0}^{t}\mathcal{A}\,\mathrm{d}t'\big]_{ij}\,
           u_j(y,t_0),
  \label{eq:matrixexp}
\end{equation}
a matrix that is independent of $y$. It maps continuous (respectively,
$C^\infty$) profiles to continuous ($C^\infty$) profiles and introduces no
spatial discontinuity, while the molecular term $\nu\partial_y^2 u_i$ is
smoothing. Continuity is preserved, establishing (ii). The forcing
displaces no fluid along $y$, so it leaves the Lebesgue measure of the
domain unchanged; the deliberate straining of the domain is treated
separately in \S\,\ref{sec:domain-strain}.
\end{proof}

The energetic role of the forcing follows from \eqref{eq:cont-evolution}.
At fixed domain length, the production term evolves the Reynolds stress as
\begin{equation}
  \left.\frac{\mathrm{d}R_{ij}}{\mathrm{d}t}\right|_{\text{prod}}
  = \overline{(\dot u_i u_j + u_i \dot u_j)}
  = -A_{ik}R_{kj}-A_{jk}R_{ki}
  \;\equiv\; \mathcal{P}_{ij},
  \label{eq:production-rate}
\end{equation}
which is \emph{exactly} the Reynolds-stress production tensor of
\eqref{eq:Pij}; the corresponding kinetic-energy production is
$\mathcal{P}\equiv\tfrac12\mathcal{P}_{ii}=-A_{ij}R_{ij}$. The production
term thus carries no model assumption. With diffusion supplying
dissipation through \eqref{eq:diffusion}, the kinetic-energy budget of the
formulation separates cleanly into an exact production by the mean strain,
a redistribution that conserves energy (the kernels), and a dissipation
(molecular diffusion). This is the structure of the exact Reynolds-stress
equation \eqref{eq:rstransport}, and it is the reason the forcing is
admissible: it adds the one term---production---that standard ODT lacks,
without disturbing the terms it already represents.

\subsection{The RDT-consistent rapid redistribution
operator}\label{sec:rapid-kernel}

\subsubsection{Why production alone is insufficient}

Were the continuous evolution to consist of production and diffusion only
($B_{ij}=0$), the line would reproduce $\mathrm{d}R_{ij}/\mathrm{d}t=
\mathcal{P}_{ij}$, whereas the exact inviscid rapid-distortion balance is
$\mathrm{d}R_{ij}/\mathrm{d}t=\mathcal{P}_{ij}+\Pi_{ij}^{(r)}$. The
omission is not small. For turbulence that is isotropic at the onset of
straining, the classical result derived below gives
$\Pi_{ij}^{(r)}=-\tfrac35\mathcal{P}_{ij}$, so that the true initial
distortion rate is $\tfrac25\mathcal{P}_{ij}$: production acting alone
overstates the rate at which anisotropy develops by a factor of
$\tfrac52$. For sustained strong strain, neglect of the pressure response
also permits a component energy to grow without the bound that
incompressibility imposes, violating realisability. Representing
$\Pi_{ij}^{(r)}$ is therefore essential, not optional.

\subsubsection{The exact rapid pressure--strain and the closure problem}

The exact rapid pressure--strain follows from the amplitude
equation~\eqref{eq:rdt}. Writing the operator there as
$\mathrm{d}\hat{u}_i/\mathrm{d}t = M_{in}\hat{u}_n$ with
$M_{in}=-A_{in}+2(k_ik_m/k^2)A_{mn}$, and noting that the incompressible
wavevector dynamics $\dot{k}_i=-A_{ji}k_j$ preserve the measure
$\mathrm{d}\boldsymbol{k}$ (its phase-space divergence is $-A_{ii}=0$), the
spectral tensor $\Phi_{ij}=\langle \hat{u}_i^{*}\hat{u}_j\rangle$ yields
\begin{equation}
  \frac{\mathrm{d}R_{ij}}{\mathrm{d}t}
  = \int\!\big(M_{in}\Phi_{nj}+M_{jn}\Phi_{in}\big)\,
    \mathrm{d}\boldsymbol{k}
  = \mathcal{P}_{ij} + \Pi_{ij}^{(r)},
  \label{eq:dRdt-rdt}
\end{equation}
in which the production~\eqref{eq:production-rate} has been separated from
the pressure contribution,
\begin{equation}
  \Pi_{ij}^{(r)} = 2A_{mn}\big(\mathcal{R}_{ijmn}+\mathcal{R}_{jimn}\big),
  \qquad
  \mathcal{R}_{ijmn} \equiv
  \int \frac{k_i k_m}{k^2}\,\Phi_{nj}(\boldsymbol{k})\,
  \mathrm{d}\boldsymbol{k}.
  \label{eq:exact-rps}
\end{equation}
The tensor $\Pi_{ij}^{(r)}$ is traceless: using the solenoidality
$k_n\Phi_{nj}=0$ (a consequence of $k_i\hat{u}_i=0$),
\begin{equation}
  \Pi_{ii}^{(r)} = 4A_{mn}\!\int\frac{k_i k_m}{k^2}\Phi_{ni}\,
  \mathrm{d}\boldsymbol{k}
  = 4A_{mn}\!\int\frac{k_m}{k^2}\big(k_i\Phi_{ni}\big)\,
  \mathrm{d}\boldsymbol{k}=0,
  \label{eq:traceless}
\end{equation}
so the rapid pressure--strain redistributes energy among components
without changing the total, in keeping with the energy-conserving
character of the ODT kernel mechanism.

Equation~\eqref{eq:exact-rps} exposes the central difficulty. The rapid
pressure--strain is a functional of the spectral tensor $\Phi_{nj}
(\boldsymbol{k})$ through the orientation-weighted integral
$\mathcal{R}_{ijmn}$; it is \emph{not} expressible in terms of the
single-point Reynolds stress $R_{ij}$ alone. This is the well-known closure
problem of the rapid pressure--strain, common to every single-point model
\citep{Crow1968,LaunderReeceRodi1975,Pope2000}; it is intrinsic to the
physics and not a deficiency peculiar to the present approach. A
one-dimensional line carries $R_{ij}$ and the one-dimensional spectrum
$\Phi_{ij}(k_2)$ along its own coordinate, but it does not carry the
distribution of energy over wavevector \emph{orientation} that
$\mathcal{R}_{ijmn}$ requires. A spectral-tensor closure is therefore
unavoidable, exactly as in second-moment closure; we make it explicit.

\subsubsection{Baseline isotropic closure}

The minimal closure assumes the spectral tensor to be instantaneously
isotropic,
$\Phi_{nj}(\boldsymbol{k})=\frac{E(k)}{4\pi k^2}\big(\delta_{nj}
-k_nk_j/k^2\big)$ with $\int E(k)\,\mathrm{d}k=k_t$. Writing
$e_i\equiv k_i/k$ and using the angular averages over the unit sphere
$\overline{e_i e_m}=\tfrac13\delta_{im}$ and
$\overline{e_i e_m e_n e_j}=\tfrac1{15}(\delta_{im}\delta_{nj}
+\delta_{in}\delta_{mj}+\delta_{ij}\delta_{mn})$,
\begin{equation}
  \mathcal{R}_{ijmn}^{\,\mathrm{iso}}
  = k_t\Big[\tfrac{4}{15}\delta_{im}\delta_{nj}
    -\tfrac1{15}\big(\delta_{in}\delta_{mj}+\delta_{ij}\delta_{mn}\big)\Big].
  \label{eq:iso-closure}
\end{equation}
Substituting \eqref{eq:iso-closure} into \eqref{eq:exact-rps} and using
$A_{nn}=0$ gives, with $S_{ij}\equiv\tfrac12(A_{ij}+A_{ji})$,
\begin{equation}
  \Pi_{ij}^{(r),\mathrm{iso}}
  = 2k_t\Big[\tfrac{4}{15}A_{ij}-\tfrac1{15}A_{ji}\Big]
   +2k_t\Big[\tfrac{4}{15}A_{ji}-\tfrac1{15}A_{ij}\Big]
  = \tfrac{4}{5}\,k_t\,S_{ij}.
  \label{eq:Pir-iso}
\end{equation}
For isotropic $R_{ij}=\tfrac23 k_t\delta_{ij}$ the production
\eqref{eq:production-rate} is $\mathcal{P}_{ij}=-\tfrac43 k_t S_{ij}$, so
that \eqref{eq:Pir-iso} is equivalent to
\begin{equation}
  \Pi_{ij}^{(r),\mathrm{iso}} = -\tfrac{3}{5}\,\mathcal{P}_{ij},
  \label{eq:crow}
\end{equation}
the exact response of initially isotropic turbulence to a weak rapid strain
\citep{Crow1968}. Equation~\eqref{eq:crow} is the content of the
isotropisation-of-production model with the constant $C_2=\tfrac35$
\citep{Pope2000}; here it is obtained, with no adjustable constant, as the
isotropic limit of the rapid-distortion integral~\eqref{eq:exact-rps}.

When the turbulence is anisotropic, \eqref{eq:iso-closure} no longer holds
and $\mathcal{R}_{ijmn}$ must be modelled as a tensor function of $R_{ij}$;
the general representation linear in the anisotropy is that of
\citet{LaunderReeceRodi1975}, which introduces calibrated constants and to
which the present construction reduces in form. We adopt
\eqref{eq:iso-closure} as the baseline and treat the anisotropic
correction as a modelling option, returning to its assessment in
\S\,\ref{sec:validation}.

\subsubsection{Realisation on the line}

Given the closed rapid pressure--strain $\Pi_{ij}^{(r)}$---symmetric and
traceless by \eqref{eq:traceless}---we require a pointwise operator
$B_{ij}$ on the line whose action reproduces it. Demanding that the
operator evolve the Reynolds stress by exactly $\Pi_{ij}^{(r)}$,
\begin{equation}
  B_{ik}R_{kj}+B_{jk}R_{ki}
  = \big(\mathsf{B}\mathsf{R}+\mathsf{R}\mathsf{B}\big)_{ij}
  = \Pi_{ij}^{(r)},
  \label{eq:lyapunov}
\end{equation}
where the second equality holds for symmetric $\mathsf{B}$. Equation
\eqref{eq:lyapunov} is a continuous Lyapunov equation; for a
positive-definite Reynolds stress $\mathsf{R}\succ0$ it possesses a unique
symmetric solution $\mathsf{B}$. The solution conserves kinetic energy
automatically: contracting \eqref{eq:lyapunov} gives
$2B_{ij}R_{ij}=\Pi_{ii}^{(r)}=0$, so that the rate of working of the
operator on the line, $\overline{B_{ij}u_iu_j}=B_{ij}R_{ij}$, vanishes. In
the isotropic case $\mathsf{R}=\tfrac23 k_t\mathsf{I}$, \eqref{eq:lyapunov}
together with \eqref{eq:Pir-iso} reduces to $\tfrac43 k_t\mathsf{B}
=\tfrac45 k_t\mathsf{S}$, that is, the closed form
\begin{equation}
  B_{ij} = \tfrac{3}{5}\,S_{ij},
  \qquad
  \mathcal{A}_{ij} = -A_{ij} + \tfrac{3}{5}\,S_{ij}.
  \label{eq:iso-B}
\end{equation}
The operator $B_{ij}$ is evaluated from the running Reynolds stress; it is
applied continuously, paced by the mean strain, exactly as is the
production. This is deliberate: the rapid pressure--strain is, by
construction, the part of the redistribution driven by the mean velocity
gradient and therefore acts at the strain rate, in contrast to the slow,
turbulence-driven redistribution that the eddy kernel
\eqref{eq:redistribution}--\eqref{eq:ci} supplies at the eddy rate. The two
mechanisms are thereby cleanly separated: strain-driven effects (production
and rapid pressure--strain) are continuous and enter through the operator
$\mathcal{A}_{ij}$, while turbulence-driven effects (the inertial cascade
and the slow return to isotropy) remain the discrete events. The uniform,
scale-independent action of $B_{ij}$ is the faithful single-point reduction
of the rapid pressure response, which in \eqref{eq:rdt} acts on every
Fourier mode through the projection $P_{ij}(\boldsymbol{k})$; the
distribution of energy across scales continues to be governed by diffusion,
the triplet-map cascade and the domain straining of
\S\,\ref{sec:domain-strain}.

We can now state precisely the sense in which the formulation is
rapid-distortion consistent.

\begin{proposition}[Rapid-distortion consistency]\label{prop:rdt}
Let $B_{ij}$ be defined by the Lyapunov equation~\eqref{eq:lyapunov} with
$\Pi_{ij}^{(r)}$ evaluated from a prescribed spectral-tensor closure. Then
the continuous evolution~\eqref{eq:cont-evolution} reproduces the
homogeneous rapid-distortion evolution of the Reynolds
stress~\eqref{eq:dRdt-rdt} exactly for that closure. With the isotropic
closure~\eqref{eq:iso-closure}, the reproduction is exact at the onset of
distortion of initially isotropic turbulence and contains no adjustable
constant, recovering the classical result~\eqref{eq:crow}.
\end{proposition}

\begin{proof}
By \eqref{eq:production-rate} the production term of
\eqref{eq:cont-evolution} contributes $\mathcal{P}_{ij}$ to
$\mathrm{d}R_{ij}/\mathrm{d}t$, and by construction~\eqref{eq:lyapunov} the
operator $B_{ij}$ contributes $\Pi_{ij}^{(r)}$. Their sum is the
right-hand side of \eqref{eq:dRdt-rdt}. For initially isotropic turbulence
$\Pi_{ij}^{(r)}$ equals the exact value~\eqref{eq:crow} by
\eqref{eq:iso-closure}--\eqref{eq:Pir-iso}.
\end{proof}

The scope of Proposition~\ref{prop:rdt} must be stated as carefully as the
result. The consistency is established at the level of the
\emph{component energies} $R_{ij}$ and is \emph{conditional on the spectral
closure}: the isotropic closure is exact only at the onset of distortion,
and its accuracy for strongly anisotropic states is that of the
isotropisation-of-production approximation it reproduces. The
\emph{spectral} distribution of the distorted energy is not enforced by
\eqref{eq:cont-evolution}; it emerges from the interaction of the strain
operator with diffusion, the cascade and the domain straining, and is
therefore a genuine prediction of the model, to be tested rather than
assumed (\S\,\ref{sec:validation}). We make no claim that the formulation
reproduces the full spectral evolution of rapid-distortion theory, which no
single-point model can; the claim is the weaker, provable one of
Proposition~\ref{prop:rdt}.

\subsection{Wavevector kinematics by domain straining}\label{sec:domain-strain}

The second ingredient of \eqref{eq:rdt}, the wavevector evolution
$\mathrm{d}k_i/\mathrm{d}t=-A_{ji}k_j$, carries the kinematic compression
and stretching of scales by the mean strain and is independent of the
amplitude evolution; in the method-of-characteristics reading of
rapid-distortion theory, the wavevectors are the characteristics along
which the amplitudes evolve. On the line, aligned with the coordinate
$x_2$, the relevant wavenumber is $k_2$. We restrict attention to the case
in which the line is aligned with a principal axis of the mean strain, so
that $A_{12}=A_{32}=0$ and a material element of the line remains a
coordinate line; this holds, in particular, for the irrotational mean flow
of a stagnation region, for which $A_{ij}$ is symmetric and its
eigenvectors define such axes. Then
\begin{equation}
  \frac{\mathrm{d}k_2}{\mathrm{d}t}=-A_{22}\,k_2 .
  \label{eq:wavenumber-line}
\end{equation}
A material element of the line of length $\ell$ stretches with the mean
flow as $\dot{\ell}/\ell=A_{22}$; since wavenumbers scale as $k_2\propto
1/\ell$, a uniform dilatation of the line reproduces
\eqref{eq:wavenumber-line} exactly,
\begin{equation}
  \frac{1}{L}\frac{\mathrm{d}L}{\mathrm{d}t}=A_{22},
  \qquad\Longrightarrow\qquad
  k_2(t)=k_2(t_0)\,\exp\!\Big(\!-\!\int_{t_0}^{t}A_{22}\,\mathrm{d}t'\Big),
  \label{eq:dilatation}
\end{equation}
where $L$ is the domain length. The dilatation~\eqref{eq:dilatation} is
implemented directly within the Lagrangian adaptive-mesh advancement of
ODT \citep{Lignell2013}, in which cell sizes already evolve, by imposing
the additional mean-flow stretching on the mesh.

This step is not redundant with the amplitude operator of
\S\S\,\ref{sec:strain-forcing}--\ref{sec:rapid-kernel}. Because the
Reynolds stress $R_{ij}$ is an intensive (per-unit-length) second moment,
it is invariant under the uniform dilatation~\eqref{eq:dilatation}: the
dilatation relabels the scales that carry the energy without changing the
component energies, while the operator $\mathcal{A}_{ij}$ changes the
component energies without relabelling scales. The two operations are the
physical-space counterparts of the characteristic motion $\boldsymbol{k}(t)$
and the amplitude evolution $\hat{u}(\boldsymbol{k}(t),t)$ in
\eqref{eq:rdt}, and together they account for the distinct contributions
of mean straining without double counting.

Two limitations are inherent to representing a three-dimensional kinematics
on one axis. The transverse wavenumbers $k_1$ and $k_3$ also strain under
$\mathrm{d}k_i/\mathrm{d}t=-A_{ji}k_j$, but are not represented on the line;
their effect on the redistribution is subsumed in the spectral closure of
\S\,\ref{sec:rapid-kernel}, which already supplies the orientation
information the line lacks. A mean flow with a rotational part, or a line
not aligned with a principal axis, tilts the material line out of the
coordinate direction and couples it to the unresolved directions; this lies
outside the present formulation and is admissible only to the extent that
the stagnation-region mean flow is irrotational.

\subsection{Summary of the formulation}\label{sec:summary}

The strain-coupled model advances the line field $u_i(y,t)$ over a
sub-step by composing three operations:
\begin{enumerate}
\item[(a)] \emph{continuous strain--diffusion phase} on a dilating domain:
  integrate $\partial_t u_i=\nu\partial_y^2 u_i+\mathcal{A}_{ij}u_j$ with
  $\mathcal{A}_{ij}=-A_{ij}+B_{ij}$, the production term exact and the
  rapid term $B_{ij}$ given by \eqref{eq:lyapunov} from the current
  Reynolds stress and the spectral closure (the baseline being
  $B_{ij}=\tfrac35 S_{ij}$, \eqref{eq:iso-B}); simultaneously dilate the
  domain by \eqref{eq:dilatation};
\item[(b)] \emph{discrete eddy events} sampled from the rate
  distribution~\eqref{eq:rate}: the triplet map~\eqref{eq:triplet} and the
  slow kernel~\eqref{eq:redistribution}--\eqref{eq:ci}, unchanged from
  standard ODT.
\end{enumerate}
Step (a) carries the mean-strain physics---production, rapid
pressure--strain and the kinematic straining of scales---and step (b)
carries the turbulence-driven physics---the inertial cascade and the slow
return to isotropy. Standard ODT is recovered exactly when $A_{ij}=0$.

The relative importance of the two steps is governed by the ratio of the
mean-strain rate $S\equiv(2S_{ij}S_{ij})^{1/2}$ to the turbulence
relaxation rate $\varepsilon/k_t$,
\begin{equation}
  \chi \equiv \frac{S\,k_t}{\varepsilon}.
  \label{eq:chi}
\end{equation}
For $\chi\gg1$ the distortion accumulates faster than the turbulence can
respond, the eddy events are negligible over the distortion time, and the
formulation reduces to the rapid-distortion limit of step (a), exact under
the closure by Proposition~\ref{prop:rdt}. For $\chi\ll1$ the turbulence
relaxes between increments of strain, the slow kernel dominates the
redistribution, and the model returns to a near-equilibrium anisotropy. The
stagnation region of a leading edge is the intermediate regime, $\chi=
O(1)$, in which neither limit is uniformly valid and both mechanisms
contribute; it is precisely there that a formulation carrying the two
together, rather than either alone, is required. The quantitative behaviour
across this range, and the fidelity of the emergent distorted spectrum, are
assessed against rapid-distortion solutions and scale-resolving data in
\S\,\ref{sec:validation}.

% =====================================================================
%  Section 4  --  Verification and validation
%  (i) analytical rapid-distortion benchmarks (derived here);
%  (ii) verification/validation protocol, metrics, failure modes.
%
%  NOTE: computed model-vs-reference comparisons (figures) are produced
%  once the simulations are run; figure environments below are clearly
%  marked placeholders and contain NO fabricated results. The analytical
%  benchmark derivations are exact and final.
%
%  Target: jfm.cls; requires amsmath. Citations resolve against references.bib.
% =====================================================================

\section{Verification and validation}\label{sec:validation}

We separate two questions. \emph{Verification} asks whether the
formulation reproduces a known exact solution in the regime where one
exists; the only such regime is the rapid-distortion limit, against which
\S\,\ref{sec:verification} tests the formulation analytically. \emph{Validation}
asks whether the formulation reproduces the behaviour of the true flow
outside that limit; \S\,\ref{sec:val-spectrum} and \S\,\ref{sec:val-les}
set out the comparisons against scale-resolving data, including the
distortion data of the present project, together with the metrics and the
diagnostic value of each possible failure. The analytical benchmarks of
\S\,\ref{sec:verification} are exact and are stated in full; the computed
comparisons are reported once the corresponding simulations are completed.

\subsection{Verification against rapid-distortion theory}\label{sec:verification}

\subsubsection{Canonical strains}

Two irrotational homogeneous strains both admit exact rapid-distortion
solutions and isolate the physics of interest. \emph{Plane strain},
\begin{equation}
  \mathsf{A}=\operatorname{diag}(a,-a,0),\qquad a>0,
  \label{eq:plane-strain}
\end{equation}
is the strain of a two-dimensional stagnation region: with $x_1$ tangential
to the surface, $x_2$ normal to it and $x_3$ spanwise, the mean flow
accelerates away from the stagnation line along $x_1$ and decelerates
towards the surface along $x_2$. The component compressed by the strain is
the wall-normal (upwash) component $u_2$, whose amplification is the
pre-impact distortion relevant to leading-edge noise. \emph{Axisymmetric
strain},
\begin{equation}
  \mathsf{A}=\operatorname{diag}(-\tfrac12\gamma,-\tfrac12\gamma,\gamma),
  \qquad \gamma>0,
  \label{eq:axi-strain}
\end{equation}
is the strain of a contraction or a three-dimensional stagnation flow and
provides an independent test in which two components are amplified and one
suppressed. For both, $\mathsf{A}$ is symmetric, so $\mathsf{A}=\mathsf{S}$
and the principal-axis alignment assumed in \S\,\ref{sec:domain-strain}
holds; the ODT line is taken along $x_2$.

\subsubsection{Onset behaviour: an exact, constant-free check}

For turbulence isotropic at the onset of straining, $R_{ij}=\tfrac23
k_t\delta_{ij}$, the continuous evolution~\eqref{eq:cont-evolution} gives,
by Proposition~\ref{prop:rdt} and equations
\eqref{eq:production-rate} and \eqref{eq:crow},
\begin{equation}
  \left.\frac{\mathrm{d}R_{ij}}{\mathrm{d}t}\right|_{0}
  = \mathcal{P}_{ij}+\Pi_{ij}^{(r),\mathrm{iso}}
  = \tfrac{2}{5}\,\mathcal{P}_{ij}
  = -\tfrac{8}{15}\,k_t\,S_{ij},
  \label{eq:onset-rate}
\end{equation}
which is the exact rapid-distortion rate for initially isotropic turbulence
and contains no adjustable constant. Equation~\eqref{eq:onset-rate} is the
sharpest verification available: it is a closed-form prediction that the
model must satisfy identically, not approximately. Evaluating it for the
two canonical strains gives the component rates in
table~\ref{tab:onset}. The wall-normal (compressed) component is amplified
in both cases; for plane strain the model predicts the upwash growth rate
\begin{equation}
  \left.\frac{\mathrm{d}R_{22}}{\mathrm{d}t}\right|_{0}
  = +\tfrac{8}{15}\,k_t\,a ,
  \label{eq:upwash-onset}
\end{equation}
to be compared with the value $\tfrac43 k_t a$ that production alone
would give: neglect of the rapid pressure--strain
overstates the upwash amplification rate by the factor $\tfrac52$
anticipated in \S\,\ref{sec:rapid-kernel}. A model that passes
\eqref{eq:onset-rate} for both strains has its production and rapid
pressure--strain terms verified jointly; failure would indicate an error
in the operator $\mathcal{A}_{ij}$ or its implementation, isolated from any
closure or cascade effect.

\begin{table}
  \begin{center}
  \def~{\hphantom{0}}
  \begin{tabular}{lccc}
    \hline
    & $\mathrm{d}R_{11}/\mathrm{d}t|_0$
    & $\mathrm{d}R_{22}/\mathrm{d}t|_0$
    & $\mathrm{d}R_{33}/\mathrm{d}t|_0$ \\[3pt]
    \hline
    Plane strain \eqref{eq:plane-strain}
      & $-\tfrac{8}{15}k_t a$ & $+\tfrac{8}{15}k_t a$ & $0$ \\[3pt]
    Axisymmetric \eqref{eq:axi-strain}
      & $+\tfrac{4}{15}k_t\gamma$ & $+\tfrac{4}{15}k_t\gamma$
      & $-\tfrac{8}{15}k_t\gamma$ \\[3pt]
    \hline
  \end{tabular}
  \caption{Exact onset rates of the component energies from
  \eqref{eq:onset-rate} for initially isotropic turbulence. The component
  compressed by the strain ($u_2$ for plane strain; $u_1,u_2$ for
  axisymmetric contraction) is amplified; the stretched component is
  suppressed.}
  \label{tab:onset}
  \end{center}
\end{table}

% =====================================================================
%  Section 4.1.3 (REVISED) + 4.1.4 (NEW)
%  Drop-in replacement for the old "Finite total strain" subsubsection
%  (which contained the [To be computed] placeholder figure fig:verification).
%
%  Figure files required (place in your figures path):
%     fig_components_plane.pdf
%     fig_components_axisymmetric.pdf
%     fig_upwash_amplification.pdf
%
%  Requires: \usepackage{graphicx}, amsmath. Citations: references.bib.
% =====================================================================

\subsubsection{Finite total strain}\label{sec:finite-strain}

Beyond onset the turbulence becomes anisotropic and the rapid
pressure--strain~\eqref{eq:exact-rps} ceases to be a function of the
Reynolds stress alone, so the verification becomes a test of the spectral
closure rather than of the construction. We integrate the model moment
equation $\mathrm{d}R_{ij}/\mathrm{d}t=\mathcal{P}_{ij}+\Pi_{ij}^{(r)}$
along each canonical strain with two closures for $\Pi_{ij}^{(r)}$: the
baseline isotropic closure~\eqref{eq:iso-closure}, equivalent to the
isotropisation-of-production model with $C_2=3/5$ (hereafter IP), and the
quasi-isotropic model of \citet{LaunderReeceRodi1975} (hereafter LRR),
\begin{equation}
  \Pi_{ij}^{(r)} = C_2\,k_t S_{ij}
  + C_3\,k_t\big(b_{ik}S_{jk}+b_{jk}S_{ik}-\tfrac{2}{3}b_{mn}S_{mn}\delta_{ij}\big)
  + C_4\,k_t\big(b_{ik}W_{jk}+b_{jk}W_{ik}\big),
  \label{eq:lrr}
\end{equation}
with $b_{ij}=R_{ij}/2k_t-\tfrac13\delta_{ij}$, mean rotation
$W_{ij}=\tfrac12(A_{ij}-A_{ji})$, and constants $C_2=4/5$, $C_3=7/4$,
$C_4=131/100$. The coefficient $C_2=4/5$ is fixed by the requirement that
\eqref{eq:lrr} reduce at isotropy ($b_{ij}=0$) to the exact value
$\tfrac45 k_t S_{ij}$ of \eqref{eq:Pir-iso}, so that LRR shares the
constant-free onset slope~\eqref{eq:onset-rate}; the strains considered
here are irrotational ($W_{ij}=0$). The benchmark is the exact
rapid-distortion solution, obtained not from a tabulated closed form but by
integrating the spectral equations~\eqref{eq:rdt} over an ensemble of
initial wavevector directions on the unit sphere (Gauss--Legendre in
$\cos\theta$ by uniform azimuth); for initially isotropic turbulence the
modal dynamics depend only on direction, so the Reynolds stress is the
spherical average of the evolved modal covariance. This integration is
independent of the moment models and reproduces the analytic onset
slope~\eqref{eq:onset-rate} to within $1.5\times10^{-3}$ (the quadrature
error), which serves as a check on the rapid-distortion machinery itself.

Figure~\ref{fig:rdt-components} reports the resulting component-energy
fractions against total strain $e=\int S\,\mathrm{d}t$ for the two strains.
Three observations follow. First, at onset all three curves for each
component are tangent: the constant-free identity~\eqref{eq:onset-rate} is
satisfied, and through $e\lesssim1$ the closures are nearly
indistinguishable from the exact solution. This is the numerical
counterpart of Proposition~\ref{prop:rdt}. Second, as anisotropy
accumulates the closures depart from the exact solution, and the LRR
closure is uniformly closer than IP: for axisymmetric strain at $e=4$ the
lateral fraction is $0.498$ (exact), $0.468$ (LRR) and $0.450$ (IP), and
the axial fraction $0.003$, $0.064$ and $0.100$ respectively. This gap is
the spectral-closure error in isolation---the eddy events are inactive in
this rapid limit---and it bounds the accuracy attainable from a single
rapid operator. Third, and most important to state plainly, for plane
strain both closures miss a \emph{qualitative} feature: the exact solution
has the spanwise component $\overline{u_3^2}$, which is neither stretched
nor compressed, grow monotonically and overtake the upwash component
near $e\approx3.3$, becoming the largest component by $e=4$
($\overline{u_3^2}/2k_t=0.51$ against $\overline{u_2^2}/2k_t=0.47$),
whereas both closures predict $\overline{u_3^2}$ to \emph{decay}. This
redistribution into the neutral direction is governed by the orientation
distribution of the spectral tensor and cannot be recovered from
single-point quantities; it is the moment-level signature of precisely the
transverse-wavenumber information that a single line does not carry
(\S\,\ref{sec:domain-strain}). It marks the total strain beyond which the
rapid closure is only qualitatively reliable, and it is one of the effects
the scale-resolving validation of \S\,\ref{sec:val-les} is designed to
probe.

\begin{figure}
  \centering
  \includegraphics[width=0.49\textwidth]{Figures/fig_components_plane.pdf}
  \hfill
  \includegraphics[width=0.49\textwidth]{Figures/fig_components_axisymmetric.pdf}
  \caption{Component energy fractions $\overline{u_i^2}/2k_t$ against total
  strain $e=St$ for initially isotropic turbulence under (\textit{a}) plane
  strain $\mathsf{A}=\mathrm{diag}(a,-a,0)$ and (\textit{b}) axisymmetric
  strain $\mathsf{A}=\mathrm{diag}(-\tfrac{\gamma}{2},-\tfrac{\gamma}{2},\gamma)$,
  normalised to $S=1$. Solid: exact rapid-distortion theory
  (wavevector-ensemble integration of \eqref{eq:rdt}); dashed: present model
  with the LRR closure~\eqref{eq:lrr}; dotted: present model with the IP
  closure~\eqref{eq:iso-closure}. Colour denotes the velocity component. For
  axisymmetric strain the two lateral components coincide,
  $\overline{u_1^2}=\overline{u_2^2}$. All closures share the exact onset
  slope; the LRR closure tracks the exact solution more closely at finite
  strain.}
  \label{fig:rdt-components}
\end{figure}

\subsubsection{Upwash amplification and the necessity of the rapid
pressure--strain}\label{sec:upwash-verif}

The component that enters the leading-edge response is the wall-normal
(upwash) component $\overline{u_2^2}$, whose spectrum is the input to
Amiet's theory. Plane strain is the relevant case, since it represents the
two-dimensional stagnation region with $x_2$ normal to the surface.
Figure~\ref{fig:rdt-upwash} isolates the evolution of the upwash fraction
under plane strain and adds, for contrast, the result of retaining
production alone ($B_{ij}=0$).

The figure makes quantitative the argument of \S\,\ref{sec:rapid-kernel}.
Production acting alone drives the upwash fraction from its isotropic value
$1/3$ to $0.98$ by $e=4$---almost the entire turbulent kinetic energy
forced into a single component, an essentially non-realisable state---and
overstates the initial amplification rate by the factor $\tfrac52$ derived
in \S\,\ref{sec:rapid-kernel}. The rapid pressure--strain removes three
fifths of this at onset~\eqref{eq:crow} and keeps the upwash fraction
physical: the exact solution rises to a peak of $0.49$ near $e\approx2.4$
and then relaxes as energy passes to the spanwise component. Both closures
reproduce the amplification of the upwash component---the very effect that
the frozen, isotropic input of classical leading-edge-noise prediction
omits---and so capture the qualitative physics that motivates the
formulation; quantitatively they overpredict the amplification at large
strain, the LRR closure ($0.60$ at $e=4$) less than IP ($0.65$), against
the exact $0.47$.

Two conclusions for the application follow. The rapid pressure--strain is
not a refinement but a necessity: without it the acoustically critical
input is wrong by a factor of order unity and unbounded in strain. And the
quantitative fidelity of the distorted upwash spectrum supplied to the
acoustic stage is set by the spectral closure and degrades with total
strain, so the accuracy of the method for a given configuration is
controlled by the strain $e=\int S\,\mathrm{d}t$ accumulated along the
stagnation streamline; for the moderate strains typical of the approach to
a leading edge the model is quantitatively reliable, while for strongly
distorted inflow the LRR closure is to be preferred. The streamwise
integral of the mean strain is therefore the natural quantity to extract
from the reference data before the model is applied, and the comparison of
the distorted upwash spectrum against scale-resolving data is taken up in
\S\,\ref{sec:val-spectrum}--\S\,\ref{sec:val-les}.

\begin{figure}
  \centering
  \includegraphics[width=0.62\textwidth]{Figures/fig_upwash_amplification.pdf}
  \caption{Upwash energy fraction $\overline{u_2^2}/2k_t$ against total
  strain under plane strain, for initially isotropic turbulence. Black:
  exact rapid-distortion theory; blue dashed: present model with the LRR
  closure; red dash-dotted: present model with the IP closure; grey dotted:
  production alone ($B_{ij}=0$). Production without the rapid
  pressure--strain overstates the onset amplification rate by $\tfrac52$ and
  forces almost all the energy into the upwash component; the rapid term
  restores the physical, bounded amplification.}
  \label{fig:rdt-upwash}
\end{figure}

% =====================================================================
%  Section 4.1.5 (REVISED) -- verification of the strain-coupled ODT
%  implementation, now with the TKE-amplification (magnitude) check.
%
%  Figure files required (place in your figures path):
%     fig_level1a_trajectory.pdf       (panel a: component fractions)
%     fig_level1a_trajectory_kt.pdf    (panel b: TKE amplification)
%
%  Requires graphicx, amsmath. Citations resolve against references.bib.
% =====================================================================

\subsubsection{Verification of the strain-coupled ODT (revise)
implementation}\label{sec:level1a}

The preceding checks verify the \emph{formulation}: that the moment
equation closed with the rapid operator reproduces rapid-distortion theory
at onset (\S\,\ref{sec:onset}) and quantify its departure at finite strain
(\S\,\ref{sec:finite-strain}). They do not test the \emph{implementation}
in the one-dimensional turbulence solver, in which the rapid operator
$B_{ij}$ is not prescribed but reconstructed at every substep: the line
Reynolds stress is formed from the resolved field, the closure supplies
$\Pi^{(r)}_{ij}$, the Lyapunov equation~\eqref{eq:lyapunov} is solved for
$B_{ij}$, and the combined operator
\[
  \mathcal{A}_{ij}=-A_{ij}+B_{ij}
\]
is applied pointwise in the advancement. An error in any of these steps
would not appear in \S\S\,\ref{sec:onset}--\ref{sec:finite-strain}.

To isolate this chain we run the solver in the rapid-distortion limit:
plane strain~\eqref{eq:plane-strain}, eddy events suppressed and molecular
diffusion set to zero, so that the continuous evolution reduces to
\[
  \partial_t u_i = \mathcal{A}_{ij}u_j
\]
acting on every cell of the line. The initial field is statistically
isotropic with
\[
  R_{ij}=\frac{2}{3}k_t\delta_{ij},
\]
and the component energies are evaluated as length-weighted averages over
the line at each output time. With the events inactive,
Proposition~\ref{prop:rdt} requires the line's component-energy evolution to
coincide with the closed moment equation; the test is whether the solver
delivers this across the whole strain range, not merely at onset.
\begin{figure}[!ht]
  \centering
  \includegraphics[width=0.49\textwidth]{Figures/fig_level1a_trajectory.pdf}
  \hfill
  \includegraphics[width=0.49\textwidth]{Figures/fig_level1a_trajectory_kt.pdf}
  \caption{Verification of the strain-coupled ODT implementation under plane
  strain, with eddy events suppressed and zero molecular diffusion.
  (\textit{a}) Component energy fractions $\overline{u_i^2}/2k_t$ (colour
  denotes the component); (\textit{b}) turbulent kinetic energy amplification
  $k_t(e)/k_t(0)$. Symbols: ODT solver, length-weighted line averages at the
  output times. Dashed: the LRR closure~\eqref{eq:lrr} integrated as a moment
  equation. Solid: exact rapid-distortion theory. The solver reproduces the
  LRR closure in the fractions (to $5\times10^{-4}$) and in the energy
  amplification (to within $3\%$); the residual energy deficit at large
  strain is interpolation loss from the adaptive mesh under-resolving the
  strained fine scales, and is addressed in \S\,\ref{sec:val-spectrum}.}
  \label{fig:level1a}
\end{figure}

Figure~\ref{fig:level1a}(\textit{a}) shows the component energy fractions
$f_{ii}\equiv\overline{u_i^2}/2k_t$ (no summation). The solver output
(symbols) lies on the LRR closure trajectory (dashed) over the entire range
$0\le e\le 4$, to within $5\times10^{-4}$ in each fraction at $e=4$. This
confirms that the reconstruction of $B_{ij}$ at each substep, the operator
assembly and the advancement are correct at every state of strain, not only
in the isotropic limit where \eqref{eq:onset-rate} is exact: the per-substep
Lyapunov solve is exercised continuously as the anisotropy develops, and it
tracks the model throughout.
 
It is essential to read the comparison correctly. The solver reproduces the
\emph{LRR closure}, and it should: with the events suppressed the line
carries no information beyond the single-point statistics that the closure
itself uses, so by construction the moment evolution is that of the closure.
The offset between the symbols and the \emph{exact} rapid-distortion
solution (solid) is therefore not an implementation error but the
single-point closure error already identified in \S\,\ref{sec:finite-strain};
agreement with the exact solution here would in fact indicate a fault, since
it would mean the solver was not evolving the prescribed closure.

The fractions are normalised quantities and so test only the
\emph{anisotropy}. The turbulent kinetic energy provides an independent
check of \emph{magnitude}, and a far less discriminating one: the rapid
pressure--strain is traceless and cannot change $k_t$ directly, which
evolves only through the production,
\begin{equation}
  \frac{\mathrm{d}\ln k_t}{\mathrm{d}e} = f_{22} - f_{11},
  \label{eq:dlnkt}
\end{equation}
integrated along the path, and that integrand is identical
across closures near onset where they agree.
Figure~\ref{fig:level1a}(\textit{b}) shows the amplification
$k_t(e)/k_t(0)$: the solver reproduces the LRR result to within $0.3\%$ over
$0\le e\le4$, confirming that the implementation captures not only the
redistribution of energy among components but its overall growth.
 
Reaching this agreement in the energy required attention to two numerical
choices that do not affect the anisotropy and are easily overlooked. First,
a broadband (white-noise) initialisation places energy at the grid scale,
where it is damped by the interpolation inherent in mesh adaption; the
resolved spectral initialisation above removes this loss, and a
mesh-refinement study confirms that the residual is independent of spatial
resolution, so it is not an under-resolution effect---consistent with the
strain operator being built from the line-averaged stress, hence spatially
uniform and generating no new scales. Second, because the inviscid limit
removes the diffusive timestep bound, the explicit advancement must be
limited instead by the imposed strain rate; left unbounded it
under-integrates the growing mode and depresses $k_t$ by a few per cent.
With a strain-resolved step both effects are removed, leaving the $0.3\%$
agreement reported above. Each acts nearly isotropically and so is invisible
in panel~(\textit{a}); the energy panel is what exposes them, and its clean
recovery completes the verification in both anisotropy and magnitude. The
resolved spectral initialisation is in any case required for the
emergent-spectrum study of \S\,\ref{sec:val-spectrum}, where the events and
the domain dilatation are activated and the spectrum---rather than the
component energies---becomes the quantity that can carry the model beyond the
single-point closure.

% =====================================================================
%  Section 4.2 — Emergent distorted spectrum  (Option A: full sweep kept)
%
%  Subsections:
%    4.2.1 linear-RDT verification (A2, eddies off + dilatation on, inviscid)
%          Figures/fig_spectrum_hs__spectra.pdf , _centroid.pdf
%    4.2.2 distorted spectrum at fixed Re (B vs C, nu=1e-5, N=100)
%          Figures/BvsC_nu1e5_spectra.pdf , _centroid.pdf
%    4.2.3 Reynolds-number dependence (full sweep, N=100/nu, under-resolved
%          low-nu cases shown WITH the resolution caveat)
%          Figures/C_ensemble_spectra.pdf , _centroid.pdf
%    4.2.4 physical operating point (nu=3e-6, N=200, resolved on 16000 cells)
%          Figures/C_kv3em6_hires_op_spectra.pdf , _centroid.pdf
%
%  Requires graphicx, amsmath, subcaption, natbib (or jfm bibliography).
%  Cross-refs out: sec:level1a, sec:gate.
% =====================================================================

\subsection{Emergent distorted spectrum}\label{sec:val-spectrum}

The verification of \S\,\ref{sec:level1a} establishes that the solver
reproduces the single-point moment evolution of rapid-distortion theory (RDT).
That is necessary but not sufficient: the single-point statistics are, by
construction, those of the closure, and a closure alone could supply them.
What ODT adds, and what Amiet's theory ultimately requires, is the
\emph{spectrum} of the distorted upwash---how the incoming turbulent energy is
redistributed across wavenumber as the gust is compressed toward the leading
edge. This section examines that spectrum: first the geometric compression in
isolation, which must reproduce linear RDT exactly; then the full model at a
fixed Reynolds number, isolated against a matched no-strain baseline; then the
dependence on Reynolds number across the resolved and under-resolved range; and
finally the physical operating point dictated by the experiment. Results are
ensemble means over independent realisations with $\pm\sigma$ bands; the
realisation count $N$ is stated for each figure.

Throughout, the temporal ODT line is mapped to the spatial leading-edge
problem through the convective hypothesis $x\approx U_\infty t$, so that the
total strain $e=St$ accumulated in time corresponds to streamwise approach
toward the stagnation region. The distortion computed here and the acoustic
response of \S\,\ref{sec:gate} are thereby treated as scale-separated: the
incoming turbulence is distorted as it convects to the leading-edge plane, and
the distorted upwash spectrum at that plane is the input to the leading-edge
response.

\subsubsection{Compression without cascade: the linear-RDT
baseline}\label{sec:spec-rdt}

With the eddy events suppressed and zero molecular viscosity, the mean strain
is the only process acting on the resolved field, and its effect is purely
kinematic. The line length contracts as
$\mathcal{L}(e)=\mathcal{L}_0\,e^{A_{22}e}$, so every resolved wavenumber is
rescaled as $k(e)=k(0)\,e^{-A_{22}e}$ with no transfer of energy between
scales: a rigid translation of the spectrum toward higher wavenumber. This is
the linear rapid-distortion result \citep{BatchelorProudman1954,HuntCarruthers1990}
for the transverse wavenumber, and it is an exact, parameter-free target.

Figure~\ref{fig:spec-rdt} confirms it. The component spectra translate toward
higher $k$ without change of shape---the peak neither broadens nor develops an
inertial range---and the spectral centroid follows $e^{-A_{22}e}$ to within
discretisation error, reaching $7.0$ at $e=3.9$ against the exact $7.03$, with
the three components migrating identically. The single-point moments are
unchanged from the no-dilatation case, since a uniform rescaling of position
leaves the length-weighted velocity statistics invariant. Compression
therefore acts entirely on the spectrum and reproduces linear RDT exactly. It
is the reference (the dashed line in all centroid figures below) against which
the full model is measured.

\begin{figure}
  \centering
  \begin{subfigure}[t]{0.48\textwidth}\centering
    \includegraphics[width=\linewidth]{Figures/fig_spectrum_hs__spectra.pdf}
    \caption{}\label{fig:rdt-spectra}
  \end{subfigure}\hfill
  \begin{subfigure}[t]{0.48\textwidth}\centering
    \includegraphics[width=\linewidth]{Figures/fig_spectrum_hs__centroid.pdf}
    \caption{}\label{fig:rdt-centroid}
  \end{subfigure}
  \caption{Compression without cascade (eddy events suppressed, inviscid).
  (\textit{a}) Component spectra $E_i(k)$ at increasing total strain (lighter to
  darker); the spectra translate to higher $k$ without changing shape.
  (\textit{b}) Spectral-centroid migration $\bar k(e)/\bar k(0)$ for each
  component (symbols) against the linear-RDT prediction $e^{-A_{22}e}$ (dashed).
  The centroid follows the prediction exactly and the three components coincide,
  confirming purely geometric compression.}
  \label{fig:spec-rdt}
\end{figure}

\subsubsection{The distorted spectrum at fixed Reynolds
number}\label{sec:spec-bvc}

The central comparison isolates what the strain does to the spectrum at a
fixed Reynolds number, holding the turbulence otherwise identical. Two
$N=100$ ensembles are run at $\nu=10^{-5}$: the full model (eddies and
dilatation active) and a baseline that is identical in every respect except
that the strain machinery is switched off, so the line does not compress.
Because the two differ only in the presence of the mean strain, the difference
between them is the distortion, cleanly separated from the cascade and
dissipation the eddies produce on their own.

Figure~\ref{fig:spec-bvc} shows the result, and it departs from linear RDT in a
specific and physically meaningful way. Consider the centroid first
(figure~\ref{fig:bvc-centroid}). The no-strain baseline, after the brief
initialisation transient near $e\simeq0.2$, relaxes monotonically to
$\bar k/\bar k_0<1$: with no compression, the eddies merely redistribute energy
and the centroid drifts below its initial value. The full model is lifted
clearly above this baseline---compression does push the spectrum toward higher
wavenumber---yet it settles at $\bar k/\bar k_0\approx1.6$--$2.0$, far below the
linear-RDT line, which rises to $7.4$ at $e=4$. The reason is visible in the
spectra (figure~\ref{fig:bvc-spectra}): the full-model spectrum is not a rigid
translation of the baseline but a \emph{broadband} distortion. It extends to
substantially higher wavenumber than the baseline---energy is genuinely carried
to small scales by compression and cascade---but the bulk of the energy remains
at the energetic scales, and what reaches high wavenumber is dissipated there
rather than accumulated. The energy-weighted centroid therefore lies well below
the value a rigid translation would give.

This is the qualitative content of the section. Linear RDT, which only
rescales the prescribed spectrum, predicts a rigid shift of the whole spectrum
to higher $k$ and a centroid on the dashed line. The full model instead
produces a broadband distorted spectrum whose shape differs from the linearly
strained one: broadened, dissipative, and with its centroid held well below the
linear prediction. The departure from linear RDT is thus real and measurable,
but it is a difference of spectral \emph{shape}---rigid translation versus
broadband cascade---not a centroid that overshoots the linear line. It is
precisely this shape difference that a frozen, linearly strained upwash
spectrum cannot represent, and that motivates supplying the ODT spectrum to
Amiet's theory in place of it, in the spirit of the rapid-distortion-modified
inflow spectra of \citet{deSantana2016} but without the linearity assumption.

\begin{figure}
  \centering
  \begin{subfigure}[t]{0.48\textwidth}\centering
    \includegraphics[width=\linewidth]{Figures/BvsC_nu1e5_spectra.pdf}
    \caption{}\label{fig:bvc-spectra}
  \end{subfigure}\hfill
  \begin{subfigure}[t]{0.48\textwidth}\centering
    \includegraphics[width=\linewidth]{Figures/BvsC_nu1e5_centroid.pdf}
    \caption{}\label{fig:bvc-centroid}
  \end{subfigure}
  \caption{Full model versus matched no-strain baseline at $\nu=10^{-5}$
  (ensemble means, $N=100$, $\pm\sigma$ band; centroid curves lightly smoothed).
  (\textit{a}) Final-strain spectrum $E(k)=\sum_i E_i(k)$: the full model
  (strain on) extends to higher wavenumber than the baseline (strain off) but
  remains broadband, not a rigid translation. (\textit{b}) Centroid migration:
  the baseline relaxes below unity, the full model is lifted above it by the
  compression, yet stays far below the linear-RDT line (dashed). The gap
  between the two solid curves is the distortion attributable to the mean
  strain at this Reynolds number.}
  \label{fig:spec-bvc}
\end{figure}

\subsubsection{Dependence on Reynolds number}\label{sec:spec-Re}

In these non-dimensional variables the molecular viscosity is the inverse
integral-scale Reynolds number, $\nu\approx Re_{\mathcal{L}}^{-1}$, and it sets
the depth of the cascade. Figure~\ref{fig:spec-Re} sweeps it over two decades
for the full model, from $\nu=10^{-4}$ to $10^{-7}$. As $\nu$ decreases the
spectrum extends progressively further toward high wavenumber
(figure~\ref{fig:C-spectra}) and the centroid rises
(figure~\ref{fig:C-centroid}): the cascade deepens monotonically with Reynolds
number, as it must. The matched no-strain baseline ensembles behave
consistently---the cascade deepens with $Re$ there too, but, with no
compression, the centroid relaxes below unity at every $\nu$, confirming that
the centroid rise in the full model is the strain's doing and not the eddies'.

Two features must be read with care, and the ensemble $\pm\sigma$ bands make
both explicit. First, the adequately resolved cases ($\nu\gtrsim10^{-5}$)
settle \emph{below} the linear-RDT line, consistent with the broadband,
dissipative distortion of \S\,\ref{sec:spec-bvc}: the eddies prevent the rigid
high-$k$ translation that inviscid linear theory predicts. Second, only the two
lowest viscosities, $\nu=10^{-6}$ and $10^{-7}$, climb toward or above the
linear line---but these are precisely the cases whose spectra flatten into a
grid-scale shelf at the highest resolved wavenumbers (figure~\ref{fig:C-spectra},
uppermost curves) and whose centroid bands are by far the widest. There the
dissipation scale has fallen below the grid, so energy piles up at the cutoff
rather than dissipating, and the centroid is inflated by an under-resolved,
partly numerical contribution. These cases are not evidence that the model
exceeds linear RDT; they are evidence that the operating point has outrun the
grid, and they fix the requirement that the chosen $\nu$ be resolution-checked
rather than simply lowered. We note also that none of the resolved cases
exhibits an extended $k^{-5/3}$ inertial range \citep{Kolmogorov1991} over this
strain interval; at these moderate Reynolds numbers the inertial range is
physically short \citep{Pope2000}, and its form is in any case resolution-limited
at the low-$\nu$ end, a matter of grid density rather than of the model.

This fixes how the Reynolds number must be chosen. It is a physical property of
the inflow turbulence the model represents---set by matching $Re_{\mathcal{L}}$
of the target inflow, hence $\nu=Re_{\mathcal{L}}^{-1}$---and not a free
parameter tuned to move the centroid relative to the linear line. Once chosen,
it must be checked for resolution: the dissipation rolloff must sit below the
grid cutoff with margin. The cases that most strongly exceed the linear line in
figure~\ref{fig:C-centroid} are exactly those failing this check, so the
physical operating point is the largest $\nu$ (lowest $Re$) consistent with the
target inflow and resolved on the available grid. That operating point is fixed
next.

\begin{figure}
  \centering
  \begin{subfigure}[t]{0.48\textwidth}\centering
    \includegraphics[width=\linewidth]{Figures/C_ensemble_spectra.pdf}
    \caption{}\label{fig:C-spectra}
  \end{subfigure}\hfill
  \begin{subfigure}[t]{0.48\textwidth}\centering
    \includegraphics[width=\linewidth]{Figures/C_ensemble_centroid.pdf}
    \caption{}\label{fig:C-centroid}
  \end{subfigure}
  \caption{Reynolds-number dependence of the full model (ensemble means,
  $N=100$ per $\nu$, $\pm\sigma$; centroid curves lightly smoothed), for
  $\nu=10^{-4},\,5\times10^{-5},\,2\times10^{-5},\,1.42\times10^{-5},\,10^{-5},
  \,10^{-6},\,10^{-7}$. (\textit{a}) Final-strain spectrum $E(k)=\sum_i E_i(k)$;
  the cascade extends further toward high $k$ as $\nu$ falls. The dotted line
  marks a $k^{-5/3}$ slope for reference. At the two smallest $\nu$ the spectrum
  flattens into a grid-scale shelf at the cutoff (under-resolved).
  (\textit{b}) Centroid migration per $\nu$ against the linear-RDT line
  (dashed). The resolved cases ($\nu\gtrsim10^{-5}$) lie below the line; the two
  lowest-$\nu$ cases exceed it but carry the widest $\sigma$ bands, reflecting
  the under-resolved grid-scale energy in (\textit{a}).}
  \label{fig:spec-Re}
\end{figure}

\subsubsection{The physical operating point}\label{sec:spec-op}

The operating point follows from \S\,\ref{sec:spec-Re}: the Reynolds number is
set by the inflow, then checked for resolution. In code units it is the inverse
integral-scale Reynolds number, $\nu\approx Re_{\mathcal{L}}^{-1}$, with
$Re_{\mathcal{L}}=u'\mathcal{L}/\nu$ built on the turbulence intensity and
integral length scale of the gust as it reaches the leading edge. For the
present configuration---a NACA0012 at $U_\infty=20$\,m\,s$^{-1}$, chord-based
Reynolds number $Re_c\approx5.1\times10^5$ and free-stream Mach number
$M=0.059$, with grid turbulence of $M=60$\,mm mesh---the turbulence at the
leading-edge plane has intensity $u'/U_\infty\approx6\%$ and integral scale
$\mathcal{L}\approx55$\,mm, giving $Re_{\mathcal{L}}\approx4.4\times10^3$.
Mapped to code units through the initial-condition normalisation
[$(u'\mathcal{L})_{\mathrm{code}}\approx1.5\times10^{-2}$, measured from the
$e=0$ field], this corresponds to $\nu\approx3\times10^{-6}$, the operating
point used here.

At this Reynolds number the dissipation scale lies below the base grid under
the sevenfold compression, so the operating point is run on a refined grid
($16000$ cells, with the minimum cell size reduced accordingly) at which the
spectrum is resolved; the grid-scale shelf of the under-resolved sweep cases in
\S\,\ref{sec:spec-Re} is absent here. The results below are ensemble means over
$N=200$ independent realisations, with bands showing the geometric $\pm\sigma$
spread across realisations.

Figure~\ref{fig:op-spectra} shows the resolved final-strain spectrum. The three
component spectra fall smoothly across four decades of wavenumber with no
grid-scale shelf, and the realisation-to-realisation spread is tight through the
energetic and inertial range, broadening only in the far dissipation tail where
it is immaterial. The transverse upwash component $E_2$ is amplified above the
streamwise component $E_1$, the spectral signature of the plane strain.
Consistent with \S\,\ref{sec:spec-Re}, no extended $k^{-5/3}$ range is present:
at $Re_{\mathcal{L}}\approx4.4\times10^3$ the inertial range is physically short
\citep{Pope2000}, and the spectrum is dominated by the energetic and dissipative
ranges. This is a property of the inflow Reynolds number, not a limitation of
the resolution, which the smooth rolloff confirms is adequate.

Figure~\ref{fig:op-centroid} shows the centroid migration. After a brief
initialisation transient near $e\simeq0.3$, in which the events rearrange the
band-limited initial field before compression dominates, the centroid settles
and migrates to $\bar k/\bar k_0\approx2.4$ by $e=4$. This is well above the
no-strain baseline (which relaxes below unity, \S\,\ref{sec:spec-bvc}) but lies
far below the linear-RDT line, which reaches $7.4$ at the same strain. The
interpretation established in \S\,\ref{sec:spec-bvc} therefore holds at the
physical operating point and with the noise averaged out over $200$
realisations: the strain-coupled model distorts the incoming spectrum into a
broadband form whose energy-weighted centroid migrates far less than the rigid
translation predicted by linear rapid-distortion theory
\citep{BatchelorProudman1954,HuntCarruthers1990}. The departure from linear
theory is robust and, at the Reynolds number of the experiment, substantial.

This resolved, broadband distorted spectrum---and in particular its upwash
component $E_2$---is the input carried forward to the leading-edge response of
\S\,\ref{sec:gate} \citep{Amiet1975}.

\begin{figure}
  \centering
  \begin{subfigure}[t]{0.48\textwidth}\centering
    \includegraphics[width=\linewidth]{Figures/C_kv3em6_hires_op_spectra.pdf}
    \caption{}\label{fig:op-spectra}
  \end{subfigure}\hfill
  \begin{subfigure}[t]{0.48\textwidth}\centering
    \includegraphics[width=\linewidth]{Figures/C_kv3em6_hires_op_centroid.pdf}
    \caption{}\label{fig:op-centroid}
  \end{subfigure}
  \caption{Physical operating point $\nu\approx3\times10^{-6}$
  ($Re_{\mathcal{L}}\approx4.4\times10^3$), resolved on a $16000$-cell grid;
  ensemble means over $N=200$ realisations with geometric $\pm\sigma$ bands.
  (\textit{a}) Final-strain component spectra $\langle E_i(k)\rangle$: the
  resolved spectrum falls smoothly across four decades with no grid-scale shelf;
  the upwash component $E_2$ (red) is amplified above the streamwise $E_1$
  (blue). The dotted line marks a $k^{-5/3}$ slope. (\textit{b}) Total-spectrum
  centroid migration $\langle\bar k(e)/\bar k(0)\rangle$ (the centroid curve
  lightly smoothed) against the linear-RDT prediction $e^{-A_{22}e}$ (dashed).
  After the initialisation transient near $e\simeq0.3$, the centroid migrates to
  $\approx2.4$ by $e=4$, far below the rigid-translation value of $7.4$:
  the broadband distortion that linear theory cannot represent.}
  \label{fig:spec-op}
\end{figure}

% =====================================================================
%  Section 5 — Closure of the rapid pressure--strain term on a single
%  ODT line (the spectral-projection question, "Gate 2").
%
%  Requires: amsmath, amssymb. Cross-refs in: sec:val-spectrum (Sec 4.2),
%  sec:level1a, eq:lrr (if defined in Sec 3). Citations: see
%  references_sec5_additions.bib (+ existing keys Crow1968,
%  LaunderReeceRodi1975, HuntCarruthers1990, Pope2000, BatchelorProudman1954).
%
%  NOTE: every modelling assumption is introduced with an explicit
%  "Assumption" label so the conditional nature of the closed form is
%  visible. The exact results (Poisson solution, M_{ijkl}) carry no
%  assumption; the single-line collapse is conditional on Assumption A1.
% =====================================================================

\section{Closure of the rapid pressure--strain term on a single line}
\label{sec:gate}

The results of \S\,\ref{sec:val-spectrum} establish the first half of the
present contribution: a gust convecting toward the leading edge is distorted
into a broadband spectrum that departs measurably---and, at the Reynolds number
of the experiment, substantially---from the rigid translation that linear
rapid-distortion theory predicts. That distortion is what the strain-coupled
ODT model is built to capture, and it is the quantity an aeroacoustic
prediction ultimately requires. Its value, however, rests on a premise that the
results alone do not justify. The rapid pressure--strain redistribution that
drives the distortion was evaluated there through the algebraic moment closure
of \S\,\ref{}%TODO-ref-Sec3
\citep{LaunderReeceRodi1975}, applied component by component on the resolved
line. That closure is a modelling input. The redistribution it represents is,
in its exact form, a \emph{nonlocal, three-dimensional} object: the rapid
pressure at a point is a volume integral over the entire fluctuating field, and
the resulting pressure--strain correlation is an integral over the full
three-dimensional velocity spectrum. Whether such an object admits any
consistent representation on the single one-dimensional line that ODT
evolves---and if so, under what condition, and at what cost---is not settled by
reproducing the moment trajectory of the closure, since the closure is by
construction able to supply that trajectory. It is a separate question, and it
is the second half of the contribution. This section answers it.

We proceed at the spectral level, where the question is sharpest. We first
write the rapid pressure--strain term exactly, with no modelling, and exhibit
the three-dimensional nonlocal integral that obstructs a single-line
representation (\S\,\ref{sec:exact-pi}--\S\,\ref{sec:obstruction}). We then
identify the one assumption under which that integral collapses to a closed
functional of the statistics the line carries, state it explicitly, justify it
for the present inflow, and bound the error it incurs
(\S\,\ref{sec:assumption}--\S\,\ref{sec:residual}). The collapsed form is a
closed, computable kernel, derived in full and verified against two exact
constraints. We show that this spectral kernel reduces, in the isotropic limit,
to precisely the moment closure used in \S\,\ref{sec:val-spectrum}
(\S\,\ref{sec:lrr-reduction})---so that the closure employed there is not an
independent assumption but the leading-order reduction of the line-consistent
term derived here. Finally, we establish that the resulting rapid source is
disjoint from the redistribution already performed by the eddy events, which is
the remaining requirement for the model to be free of double counting
(\S\,\ref{sec:doublecount}). The outcome is conditional rather than universal:
the rapid term \emph{does} close on a single line, within a stated and
physically grounded limit, and the limit is exactly the regime the
leading-edge problem occupies.

\subsection{The exact rapid pressure--strain term}\label{sec:exact-pi}

For incompressible flow the fluctuating pressure satisfies a Poisson equation
whose source splits into a part linear in the fluctuations and the mean
velocity gradient (the \emph{rapid} part) and a part quadratic in the
fluctuations (the \emph{slow} part) \citep{Pope2000}:
\begin{equation}
  \frac{1}{\rho}\,\nabla^2 p'
  = \underbrace{-2\,\frac{\partial U_k}{\partial x_l}\,
      \frac{\partial u'_l}{\partial x_k}}_{\text{rapid}}
    \;\underbrace{-\,\frac{\partial^2}{\partial x_k\partial x_l}
      \bigl(u'_k u'_l - \overline{u'_k u'_l}\bigr)}_{\text{slow}} .
  \label{eq:poisson}
\end{equation}
The rapid part is linear in $\boldsymbol{u}'$ and proportional to the mean
gradient $\partial U_k/\partial x_l$, which for the present homogeneous plane
strain is the constant tensor $A_{kl}$ of \S\,\ref{sec:val-spectrum}. Writing
$p'_{\mathrm r}$ for the rapid pressure and inverting the Laplacian with the
free-space Green's function $G(\boldsymbol{x})=-1/4\pi|\boldsymbol{x}|$,
\begin{equation}
  p'_{\mathrm r}(\boldsymbol{x})
  = \frac{\rho}{2\pi}\,A_{kl}\!\int
    \frac{1}{|\boldsymbol{x}-\boldsymbol{x}'|}\,
    \frac{\partial u'_l}{\partial x'_k}(\boldsymbol{x}')\,
    \mathrm{d}\boldsymbol{x}' .
  \label{eq:prapid}
\end{equation}
The rapid pressure--strain correlation that redistributes energy among the
Reynolds-stress components is
$\Pi^{(\mathrm r)}_{ij}
   = \overline{(p'_{\mathrm r}/\rho)\,(\partial u'_i/\partial x_j
     + \partial u'_j/\partial x_i)}$.
Substituting \eqref{eq:prapid} and using homogeneity to pass to the
velocity spectrum tensor
$\Phi_{mn}(\boldsymbol{\kappa})$
(the Fourier transform of $\overline{u'_m(\boldsymbol{x})\,
u'_n(\boldsymbol{x}+\boldsymbol{r})}$),
one obtains the standard exact result
\citep{BatchelorProudman1954,Crow1968,HuntCarruthers1990}
\begin{equation}
  \Pi^{(\mathrm r)}_{ij}
  = A_{kl}\,M_{ijkl},
  \qquad
  M_{ijkl}
  = \int_{\mathbb{R}^3}
    \Bigl[\,
      \frac{\kappa_j\kappa_k}{\kappa^2}\,\Phi_{il}(\boldsymbol{\kappa})
      + \frac{\kappa_i\kappa_k}{\kappa^2}\,\Phi_{jl}(\boldsymbol{\kappa})
    \,\Bigr]\,\mathrm{d}\boldsymbol{\kappa},
  \label{eq:Mijkl}
\end{equation}
where $\kappa=|\boldsymbol{\kappa}|$. Equations
\eqref{eq:poisson}--\eqref{eq:Mijkl} are exact; no modelling has entered. The
tensor $M_{ijkl}$ is a linear functional of the \emph{full three-dimensional}
spectrum $\Phi_{mn}(\boldsymbol{\kappa})$ through the nonlocal projection
operator $\kappa_a\kappa_b/\kappa^2$, which is the wavenumber-space image of
the inverse Laplacian in \eqref{eq:prapid}. This nonlocality and
three-dimensionality are the obstruction that the single-line model must
confront.

\subsection{The obstruction: what an ODT line resolves}\label{sec:obstruction}

ODT evolves a velocity vector $\boldsymbol{u}'(y)$ on a single line aligned
with the resolved direction, here taken along the contraction axis of the
strain, $y\equiv x_2$. From it the model has access to the one-dimensional
spectra
\begin{equation}
  \phi_{mn}(\kappa_2)
  = \int_{\mathbb{R}^2}
      \Phi_{mn}(\boldsymbol{\kappa})\,
      \mathrm{d}\kappa_1\,\mathrm{d}\kappa_3,
  \label{eq:phi1d}
\end{equation}
the one-dimensional spectra obtained by integrating $\Phi_{mn}$ over the two
unresolved wavenumber components. The integrand of \eqref{eq:Mijkl}, by
contrast, carries the \emph{direction-resolved} weight
$\kappa_a\kappa_b/\kappa^2$, which cannot be reconstructed from
$\phi_{mn}(\kappa_2)$ alone: integrating out $\kappa_1,\kappa_3$ as in
\eqref{eq:phi1d} destroys exactly the angular information that the projection
operator acts on. The reduction $M_{ijkl}\!\to\!$ (functional of
$\phi_{mn}$) is therefore not available without an assumption about how the
energy is distributed over the unresolved wavevector directions. This is the
crux of the present question, and we make the required assumption explicit
rather than implicit.

\subsection{The closure assumption and the collapsed form}
\label{sec:assumption}

\paragraph{Assumption A1 (axisymmetry of the undistorted spectrum about the
line).} The spectrum tensor of the incoming turbulence, before the rapid
distortion acts, is axisymmetric about the resolved direction $\boldsymbol{e}_2$:
$\Phi_{mn}(\boldsymbol{\kappa})$ depends on $\boldsymbol{\kappa}$ only through
$\kappa_2$ and $\kappa_\perp=(\kappa_1^2+\kappa_3^2)^{1/2}$, and its
polarization is invariant under rotation about $\boldsymbol{e}_2$.

This is a statement about the \emph{turbulence}, not about the strain. It is
justified for the present configuration because the inflow is grid turbulence,
which is closely isotropic at the grid and is subsequently contracted along a
single axis as it approaches the stagnation plane; isotropy is the special
case of A1, and uniaxial contraction preserves axisymmetry about the
contraction axis. The plane strain $A_{kl}=\mathrm{diag}(\tfrac12,-\tfrac12,0)$
acting at the leading edge is itself \emph{not} axisymmetric, so A1 is imposed
on the state on which the rapid operator acts, consistent with the
rapid-distortion ordering in which $\Pi^{(\mathrm r)}$ is evaluated on the
current (here, axisymmetric to leading order) field. The error incurred when
the accumulated distortion has driven the spectrum away from axisymmetry is
quantified in \S\,\ref{sec:residual}.

Under A1 the angular integral in \eqref{eq:Mijkl} can be carried out. Writing
$\boldsymbol{\kappa}=(\kappa_\perp\cos\theta,\ \kappa_2,\
\kappa_\perp\sin\theta)$ and using that, by A1, $\Phi_{mn}$ is independent of
$\theta$, the azimuthal averages of the projection weights reduce to
\begin{equation}
  \Bigl\langle \frac{\kappa_1^2}{\kappa^2}\Bigr\rangle_\theta
  = \Bigl\langle \frac{\kappa_3^2}{\kappa^2}\Bigr\rangle_\theta
  = \frac{1}{2}\,\frac{\kappa_\perp^2}{\kappa^2},
  \qquad
  \Bigl\langle \frac{\kappa_2^2}{\kappa^2}\Bigr\rangle_\theta
  = \frac{\kappa_2^2}{\kappa^2},
  \qquad
  \Bigl\langle \frac{\kappa_1\kappa_3}{\kappa^2}\Bigr\rangle_\theta = 0,
  \label{eq:angular}
\end{equation}
with $\kappa^2=\kappa_2^2+\kappa_\perp^2$, and all weights mixing a resolved
and an unresolved index (e.g.\ $\kappa_1\kappa_2/\kappa^2$) averaging to zero.
Substituting \eqref{eq:angular} into \eqref{eq:Mijkl} removes the azimuthal
dependence and leaves a two-dimensional integral over $(\kappa_2,\kappa_\perp)$
of the axisymmetric spectrum. Defining the diagonal one-dimensional spectra
along the line,
\begin{equation}
  E_n^{\,\|}(\kappa_2)
   = \int_0^\infty \Phi_{nn}(\kappa_2,\kappa_\perp)\,
     2\pi\kappa_\perp\,\mathrm{d}\kappa_\perp
   \quad(\text{no sum}),
   \label{eq:Epar}
\end{equation}
the rapid pressure--strain term reduces to a two-dimensional integral over the
half-plane $(\kappa_2,\kappa_\perp)$ of the axisymmetric spectrum. The
axisymmetric, solenoidal spectrum is fixed by two scalar functions of
$(\kappa_2,\kappa_\perp)$ \citep{Batchelor1946,ChandrasekharAxisym1950}; we
denote them $a(\kappa_2,\kappa_\perp)$ and $c(\kappa_2,\kappa_\perp)$, where
$a$ is the isotropic-like scalar (the sole survivor when the turbulence is
isotropic) and $c$ measures the axisymmetric departure from isotropy. Carrying
out the azimuthal average \eqref{eq:angular} and specialising to the plane
strain $A=\mathrm{diag}(\tfrac12,-\tfrac12,0)$, the diagonal rapid
pressure--strain components are
\begin{equation}
  \Pi^{(\mathrm r)}_{nn}
  = \int_{0}^{\infty}\!\!\int_{0}^{\infty}
      g_{n}(x)\,\bigl[\,a(\kappa_2,\kappa_\perp),\,c(\kappa_2,\kappa_\perp)\,\bigr]\,
      2\pi\kappa_\perp\,\mathrm{d}\kappa_\perp\,\mathrm{d}\kappa_2
   \quad(\text{no sum on }n),
  \label{eq:Piclosed}
\end{equation}
with $x\equiv\kappa_2^2/\kappa^2$, $\kappa^2=\kappa_2^2+\kappa_\perp^2$, and the
kernels, linear in the two spectral scalars,
\begin{align}
  g_{1}(x) &= a\!\left(\tfrac18 + \tfrac34 x - \tfrac78 x^2\right)
            + c\,\tfrac78\,x\,(1-x)^2, \label{eq:g1}\\[2pt]
  g_{2}(x) &= -\tfrac32\,x\,\bigl[\,a\,(1-x) + c\,(1-x)^2\,\bigr],
            \label{eq:g2}\\[2pt]
  g_{3}(x) &= a\!\left(-\tfrac18 + \tfrac34 x - \tfrac58 x^2\right)
            + c\,\tfrac58\,x\,(1-x)^2. \label{eq:g3}
\end{align}
Equations \eqref{eq:Piclosed}--\eqref{eq:g3} are the central result:
\emph{under A1, the rapid pressure--strain term is a closed, computable
functional of the axisymmetric spectral scalars $a,c$ on the resolved
half-plane and the mean strain}, with no reference to the unresolved azimuthal
structure beyond what A1 fixes. The kernels satisfy two exact constraints that
verify the reduction. First, energy conservation: pressure--strain
redistributes without changing the trace, and indeed
$g_1(x)+g_2(x)+g_3(x)=0$ identically for all $x$ and for both scalars, so
$\Pi^{(\mathrm r)}_{11}+\Pi^{(\mathrm r)}_{22}+\Pi^{(\mathrm r)}_{33}=0$.
Second, the isotropic limit: setting $c=0$ with $a$ a function of $\kappa$
alone, the integral of $g_3$ over the half-plane vanishes, so
$\Pi^{(\mathrm r)}_{33}=0$---as it must, since the plane strain has no
component along $x_3$ and the isotropic rapid term is proportional to $S_{ij}$
\citep{Crow1968}. Both constraints hold identically in the derived kernels,
confirming the tensor reduction.

\subsection{Residual error and the limit of validity}\label{sec:residual}

Assumption A1 is exact for the undistorted inflow and is progressively
violated as the rapid strain itself renders the spectrum anisotropic about the
line. The leading correction is governed by the departure-from-axisymmetry
tensor $D_{mn}(\boldsymbol{\kappa})
=\Phi_{mn}-\Phi^{\mathrm{axi}}_{mn}$, whose contribution to
\eqref{eq:Mijkl} is the azimuthal first harmonic that \eqref{eq:angular}
discards. Because that harmonic integrates to zero against the isotropic part
and enters $M_{ijkl}$ only at second order in the accumulated anisotropy, the
closed form \eqref{eq:Piclosed} is accurate to $O(b^2)$ in the Reynolds-stress
anisotropy $b_{ij}$ generated by the strain, and exact at $b_{ij}=0$. The
practical consequence is that \eqref{eq:Piclosed} is most reliable early in
the distortion, when the gust first enters the strained region, and degrades
as the total strain $e$ accumulates---precisely the regime ordering that the
convective mapping of \S\,\ref{sec:val-spectrum} makes physical, since the
acoustically relevant upwash is set near the leading-edge plane where the
accumulated strain is bounded. We therefore regard \eqref{eq:Piclosed} as a
controlled leading-order closure, valid where the gust anisotropy is moderate,
rather than as an identity.

\subsection{Reduction to the moment closure}\label{sec:lrr-reduction}

The closed form \eqref{eq:Piclosed}--\eqref{eq:g3} operates on the resolved
spectrum, whereas the distortion results of \S\,\ref{sec:val-spectrum} were
obtained with the algebraic moment closure of \S\,\ref{}%TODO-ref-Sec3
\citep{LaunderReeceRodi1975}. For the two to be consistent, the spectral
kernel must reduce to that moment closure in the limit where the moment
closure is exact---isotropic turbulence. We show that it does, which both
justifies the closure used in the simulations and fixes the prefactor in an
explicit convention.

Figure~\ref{fig:pi-kernels} plots the kernels $g_1,g_2,g_3$ of
\eqref{eq:g1}--\eqref{eq:g3} against the angular variable
$x=\kappa_2^2/\kappa^2$. Two features are exact and hold for both spectral
scalars. The kernels sum to zero at every $x$,
\begin{equation}
  g_1(x)+g_2(x)+g_3(x)=0 \qquad \forall x,
  \label{eq:gtrace}
\end{equation}
so the rapid term conserves turbulent kinetic energy pointwise in wavenumber,
not merely in the integral; and each kernel vanishes at $x=1$, where the
wavevector aligns with the line and the projection $\kappa_a\kappa_b/\kappa^2$
degenerates. The signs are physically ordered: $g_1>0$ (energy fed to the
streamwise component $E_1$), $g_2<0$ (energy removed from the line/upwash
component $E_2$) over the bulk of the angular range, with $g_3$ changing sign,
consistent with the plane strain stretching $x_1$ and compressing $x_2$.

For \emph{isotropic} turbulence the anisotropy scalar vanishes ($c=0$) and the
isotropic-like scalar $a$ depends only on $\kappa$. Averaging the kernels over
solid angle---equivalently, integrating $g_n(x)$ against the isotropic
distribution of $x$ on the sphere---gives
\begin{equation}
  \langle g_1\rangle : \langle g_2\rangle : \langle g_3\rangle
  = +\tfrac15 : -\tfrac15 : 0,
  \label{eq:giso}
\end{equation}
in the ratio $+1:-1:0$. This is exactly the structure of the plane strain
$S_{ij}=\mathrm{diag}(\tfrac12,-\tfrac12,0)$: the rapid term is proportional to
$S_{ij}$, with $\Pi^{(\mathrm r)}_{33}=0$ because the strain has no component
along $x_3$. The spectral kernel therefore collapses, under isotropy, to the
isotropic-production form of the moment closure,
\begin{equation}
  \Pi^{(\mathrm r)}_{ij}
  = \tfrac{4}{5}\,k_t\,S_{ij}
  \qquad(\text{isotropic limit}),
  \label{eq:crow}
\end{equation}
the result of \citet{Crow1968} that underlies the isotropisation-of-production
(IP) rapid term in the closure of \citet{LaunderReeceRodi1975}. We note a
convention: with the symmetrised definition of the pressure--strain correlation
adopted here and the factor of two carried in the Poisson source
\eqref{eq:poisson}, the kernels \eqref{eq:g1}--\eqref{eq:g3} yield the
coefficient $\tfrac25 k_t$ for the $11$-component directly (since $S_{11}=\tfrac12$),
equivalent to the standard $\tfrac45 k_t S_{ij}$ of \eqref{eq:crow}; the ratio
\eqref{eq:giso} and the tracelessness \eqref{eq:gtrace} are independent of this
convention.

The consequence is that the algebraic closure used to obtain the distortion of
\S\,\ref{sec:val-spectrum} is not an independent assumption but the
isotropic-limit reduction of the line-consistent spectral term derived here.
The spectral form \eqref{eq:Piclosed} carries, in addition, the angular
($x$-dependent) and anisotropy ($c$-scalar) information that the moment closure
discards; the corrections it retains are precisely those of order $b^2$ in the
strain-generated anisotropy identified in \S\,\ref{sec:residual}. The moment
closure and the spectral closure thus agree to leading order and depart in a
controlled way as the gust anisotropy grows, which is the regime ordering the
present application requires.

\begin{figure}
  \centering
  \includegraphics[width=0.72\textwidth]{Figures/fig_pi_kernels.pdf}
  \caption{Rapid pressure--strain closure kernels $g_1,g_2,g_3$
  (eqs.~\eqref{eq:g1}--\eqref{eq:g3}) against the angular variable
  $x=\kappa_2^2/\kappa^2$. Solid curves: coefficient of the isotropic-like
  spectral scalar $a$; faint dashed curves: coefficient of the anisotropy
  scalar $c$. The kernels sum to zero at every $x$ (energy conservation,
  eq.~\eqref{eq:gtrace}) and vanish at $x=1$. Averaged over an isotropic
  spectrum they give the ratio $+\tfrac15:-\tfrac15:0$
  (eq.~\eqref{eq:giso}), recovering the isotropisation-of-production form
  $\Pi^{(\mathrm r)}_{ij}=\tfrac45 k_t S_{ij}$ \citep{Crow1968} that underlies
  the moment closure used in \S\,\ref{sec:val-spectrum}.}
  \label{fig:pi-kernels}
\end{figure}

\subsection{Consistency with the eddy redistribution: avoiding double
counting}\label{sec:doublecount}

A closed rapid term is necessary but not sufficient: it must not duplicate
redistribution already effected by the eddy events. The pressure--strain
correlation decomposes canonically into rapid and slow parts,
$\Pi_{ij}=\Pi^{(\mathrm r)}_{ij}+\Pi^{(\mathrm s)}_{ij}$, corresponding to the
two source terms in \eqref{eq:poisson}. The slow part
$\Pi^{(\mathrm s)}_{ij}$ is the return-to-isotropy mechanism, quadratic in the
fluctuations and independent of the mean gradient; the eddy events in ODT,
which act on the fluctuating field without reference to $A_{kl}$, are the
model's representation of exactly this slow redistribution. The rapid part
$\Pi^{(\mathrm r)}_{ij}$, proportional to $A_{kl}$ and linear in the
fluctuations, is by construction absent from the eddy mechanism, which carries
no dependence on the mean strain. The two therefore occupy disjoint terms of
the canonical decomposition, and adding \eqref{eq:Piclosed} as an explicit
mean-strain-driven source alongside the eddies does not double-count, provided
the eddy events are calibrated to reproduce the slow return-to-isotropy in the
absence of mean strain---a condition that the no-strain baseline of
\S\,\ref{sec:spec-bvc} verifies directly, since there the centroid relaxes
under the eddies alone with no spurious rapid contribution. This argument
establishes consistency at the level of the decomposition; a quantitative
check that the calibrated slow part matches the measured return-to-isotropy is
deferred to \S\,\ref{}.%TODO-ref

\subsection{Summary}\label{sec:gate-summary}

The rapid pressure--strain term, exactly a nonlocal three-dimensional integral
\eqref{eq:Mijkl} of the full spectrum, collapses under axisymmetry of the
undistorted inflow about the resolved direction (Assumption A1) to the closed,
single-line functional \eqref{eq:Piclosed} of the resolved spectra and the
mean strain. The collapse is exact for the isotropic inflow and controlled to
$O(b^2)$ in the strain-generated anisotropy, with the error growing as the
accumulated strain renders the spectrum anisotropic about the line. The
resulting rapid source is consistent with---and disjoint from---the slow
redistribution carried by the eddy events, so the strain-coupled model is free
of double counting at the level of the rapid/slow decomposition. The closure
therefore holds on a single line within a stated and physically grounded
limit, rather than universally; characterising its breakdown at large
accumulated strain is the subject of \S\,\ref{}.%TODO-ref

% =====================================================================
%  Validation against Chen, Meneveau & Katz (2006) -- Option A:
%  rigorous RDT-reference validation + qualitative full-model anisotropy.
%  Figures: fig_cmk_rdt_validation.pdf, fig_cmk_anisotropy.pdf
%  Citations (existing keys): ChenMeneveauKatz2006, LeePiomelliReynolds1986,
%  BatchelorProudman1954. Add the first two to references if absent.
% =====================================================================

\subsection{Validation against a strained-turbulence benchmark}
\label{sec:cmk}

The verification of \S\,\ref{sec:val-spectrum} establishes that the model
departs from linear rapid-distortion theory in a broadband manner, but it does
so against the model's own no-strain baseline. An external benchmark is needed
to confirm that both the rapid-distortion limit and the departure from it are
physical. We use the plane-strain experiment of \citet{ChenMeneveauKatz2006}
(hereafter CMK), in which initially near-isotropic turbulence at
$R_\lambda\approx400$ is subjected to a controlled planar straining cycle and
the one-dimensional component spectra are compared with rapid-distortion theory.
We reproduce the \emph{straining} phase of their cycle; the relaxation and
destraining phases involve a strain reversal that the present monotonic-strain
configuration does not model. In CMK's variables the end of the straining phase
corresponds to a total deformation $D=\exp(\int S\,\mathrm{d}t)\approx3$, which
maps to a total strain $e=St=2\ln D\approx2.20$ in the present notation.

\begin{figure}
  \centering

  \begin{subfigure}[t]{0.48\textwidth}
    \centering
    \includegraphics[width=\linewidth]{Figures/fig_cmk_centroid.pdf}
    \caption{}
    \label{fig:cmk-rdt-centroid}
  \end{subfigure}
  \hfill
  \begin{subfigure}[t]{0.48\textwidth}
    \centering
    \includegraphics[width=\linewidth]{Figures/fig_cmk_energy.pdf}
    \caption{}
    \label{fig:cmk-rdt-energy}
  \end{subfigure}

  \caption{Rapid-distortion limit, with eddy events suppressed and inviscid
  evolution, against the analytic reference of \citet{LeePiomelliReynolds1986}
  used by \citet{ChenMeneveauKatz2006}. (\textit{a}) Component
  spectral-centroid migration follows the rigid-translation prediction
  $e^{-A_{22}e}$, shown by the solid line, reaching the deformation ratio
  $D=3$ at the end of the straining phase, with the three components
  coincident. (\textit{b}) Turbulent kinetic energy follows the plane-strain
  amplification $q^2/q_0^2$, shown by the solid line, reaching $1.74$ at
  $e=2.20$.}
  \label{fig:cmk-rdt}
\end{figure}
With the eddy events suppressed and zero molecular viscosity, the model must
reproduce the analytic rapid-distortion result that CMK use as their reference,
taken from \citet{LeePiomelliReynolds1986}. Figure~\ref{fig:cmk-rdt} shows the
comparison. The spectral centroid of each velocity component
(figure~\ref{fig:cmk-rdt-centroid}) follows the rigid-translation prediction
$\bar\kappa(e)/\bar\kappa(0)=e^{-A_{22}e}$ exactly, reaching the deformation
ratio $D=3.00$ at $e=2.20$, with the three components coincident. The turbulent
kinetic energy (figure~\ref{fig:cmk-rdt-energy}) follows the
\citet{LeePiomelliReynolds1986} plane-strain amplification
$q^2/q_0^2 = 1 + \tfrac{8}{15}a^2 + \tfrac{8}{315}a^4$ (with $a=\int S\,\mathrm
{d}t$), reaching $1.74$ against the analytic $1.74$ at the end of the straining
phase. The strain operator therefore reproduces the rapid-distortion limit that
underlies CMK's reference, to within discretisation error, in both the spectral
translation and the energy amplification. We note that this agreement requires
the truly inviscid setting ($\nu=0$): on a fine adaptive grid any nonzero
viscosity dissipates the compressed small scales and degrades the rigid
translation, an effect we verified directly.

\subsubsection{The full model: emergence of spectral
anisotropy}\label{sec:cmk-full}

With the eddy events and molecular viscosity active, the model is no longer
constrained to rigid translation, and the strain drives a measurable anisotropy
between the velocity components. Figure~\ref{fig:cmk-aniso} shows the
ensemble-averaged ($N=64$) centroid migration of the streamwise ($E_{11}$) and
upwash ($E_{22}$) component spectra. The band-limited initial field is not an
equilibrium turbulent state; over the first strain increment the eddy events
rearrange it to a relaxed state, accompanied by a rapid initial energy
adjustment (shaded region). Measured from this relaxed reference state, the two
components diverge progressively under continued strain: the streamwise
component migrates to higher wavenumber while the upwash component is held back.
This is the same qualitative behaviour CMK report---the strained spectrum
departs from the isotropic-translation expectation, with the components
responding differently to the compression---and it is the spectral signature of
the broadband distortion identified in \S\,\ref{sec:val-spectrum}.

We restrict this comparison to its qualitative content for three reasons, each
stated plainly. First, the present band-limited initial spectrum differs from
CMK's model spectrum, so absolute spectral levels are not directly comparable;
the relaxation transient further precludes an absolute-level match. Second, the
one-dimensional spectra resolved on the ODT line and the transverse/longitudinal
one-dimensional spectra of the three-dimensional experiment are not the same
projection of the spectrum tensor, so the cross-component ratio differs by
construction even in the isotropic state. Third, the effective scale separation
achieved here, $\mathcal{L}/\eta\approx250$, is below the experimental value
$\approx930$; matching the latter on a single line exceeds feasible cost, as the
compressed dissipation range demands prohibitive resolution. Within these
limits, the model reproduces the rapid-distortion reference exactly
(\S\,\ref{sec:cmk-rdt}) and the qualitative departure mechanism of the full
benchmark (this section). A quantitative spectral reproduction, requiring a
matched equilibrium inflow spectrum and matched Reynolds number, is left to
future work.



\begin{figure}
  \centering
  \includegraphics[width=0.7\textwidth]{Figures/fig_cmk_anisotropy.pdf}
  \caption{Full model (eddies and viscosity active; ensemble mean, $N=64$):
  centroid migration of the streamwise ($E_{11}$) and upwash ($E_{22}$)
  component spectra, each normalised to the relaxed reference state
  $e_0=0.1$. The band-limited initial field relaxes over the first strain
  increment (shaded); thereafter the components diverge under strain---the
  streamwise component migrating to higher wavenumber while the upwash component
  is held back---reproducing the qualitative departure-from-RDT mechanism of
  \citet{ChenMeneveauKatz2006}. The comparison is qualitative; see text.}
  \label{fig:cmk-aniso}
\end{figure}


\newpage
% =====================================================================
%  Validation against Lee & Reynolds (1985) DNS of strained homogeneous
%  turbulence. Moment-level (Reynolds-stress anisotropy) comparison.
%
%  Figures: fig_LR_plane.pdf, fig_LR_axisym.pdf
%  Bib keys (add via references_lee_reynolds_additions.bib):
%     LeeReynolds1985, MillsCorrsin1959, TuckerReynolds1968
%  Reuses (already in manuscript): LeePiomelliReynolds1986, BatchelorProudman1954,
%     ChenMeneveauKatz2006, Pope2000, LaunderReeceRodi1975.
% =====================================================================

\subsection{Validation against direct numerical simulation of strained
turbulence}
\label{sec:lr-validation}

The benchmark of \S\,\ref{sec:cmk} establishes that the present model reproduces
the rapid-distortion limit exactly and departs from it in the broadband manner
expected at finite Reynolds number. That comparison is, however, restricted in
two respects: it is an experiment, and at its Reynolds number
($R_\lambda\approx400$) a quantitative spectral reproduction on a single
one-dimensional line is precluded by resolution. To validate the model where it
is most directly applicable---at the level of the single-point Reynolds-stress
anisotropy, which is the quantity that ultimately feeds the leading-edge
response---we compare against the direct numerical simulations (DNS) of
\citet{LeeReynolds1985}. Their simulations are an ideal second benchmark for
three reasons. First, they are DNS, free of the measurement noise and finite-%
sample scatter of an experiment. Second, they were deliberately run at low
turbulence Reynolds number, $q^4/(\nu\varepsilon)\lesssim100$, a constraint
imposed in the original study by the requirement to resolve the small scales on
a $128^3$ mesh; this same low Reynolds number is fully resolvable on a single
ODT line, so the resolution wall encountered in \S\,\ref{sec:cmk} does not arise.
Third, they report the Reynolds-stress anisotropy as a function of total strain
for several irrotational strain modes, alongside the analytic rapid-distortion
prediction, providing a graded reference curve rather than a single endpoint.

\subsubsection{Comparison quantity and strain protocol}
\label{sec:lr-quantity}

The comparison is made in terms of the Reynolds-stress anisotropy tensor
\begin{equation}
  b_{ij} \;=\; \frac{\langle u_i u_j\rangle}{\langle u_k u_k\rangle}
              \;-\; \tfrac{1}{3}\,\delta_{ij},
  \label{eq:bij}
\end{equation}
evaluated as a function of the reference total strain
\begin{equation}
  c \;=\; \exp\!\left(\int_0^t S(t')\,\mathrm{d}t'\right),
  \qquad
  S \;=\; \sqrt{\tfrac{1}{2}\,S_{ij}S_{ij}},
  \label{eq:total-strain}
\end{equation}
where $S$ is the equivalent mean strain rate and $c$ is the cumulative
deformation, following the definitions of \citet[][Eqs.~5.1.4 and
5.1.7]{LeeReynolds1985}. Note that $c$ is a deformation \emph{ratio}: it is
related to the integrated strain in the present notation by
$\int_0^t S\,\mathrm{d}t = \ln c$. The anisotropy~\eqref{eq:bij} is computed
directly from the single-point velocity-component variances along the ODT line,
ensemble-averaged over $N=1000$ independent realizations; no spectral transform
is involved, so the comparison is insensitive to the small-scale resolution that
limited the spectral comparison of \S\,\ref{sec:cmk}.

Two of \citeauthor{LeeReynolds1985}'s irrotational strain modes are considered.
The first is \emph{plane strain},
\begin{equation}
  S_{ij} \;=\; S\,\mathrm{diag}(0,\,-1,\,+1),
  \label{eq:plane-strain}
\end{equation}
in which one direction is unstrained, one is compressed, and one is stretched.
The second is \emph{axisymmetric contraction},
\begin{equation}
  S_{ij} \;=\; \frac{2}{\sqrt{3}}\,S\,
  \mathrm{diag}\!\left(+1,\,-\tfrac{1}{2},\,-\tfrac{1}{2}\right),
  \label{eq:axisym-contraction}
\end{equation}
in which one axis is stretched and the two transverse directions are compressed
equally. In both cases the ODT line is oriented along a compressed direction, so
that the mean dilatation compresses the line as $L(t)=L_0\exp(A_{22}t)$, with
$A_{22}=-S$ for plane strain and $A_{22}=-S/\sqrt{3}$ for axisymmetric
contraction. The plane-strain runs are carried to $c=4$ and the axisymmetric runs
to $c=3.32$, matching the total strains reported by \citeauthor{LeeReynolds1985}
The kinematic fidelity of the strain implementation is confirmed by the line
compression: the measured $L(t_f)/L_0$ agrees with $\exp(A_{22}t_f)$ to better
than $0.3\%$ in both cases.

A note on the strain-rate parameter is in order, because it governs how the two
modes should be interpreted. \citeauthor{LeeReynolds1985} characterise each run
by the strain-rate parameter $S^\ast=Sq^2/\varepsilon$, the ratio of the
eddy-turnover time to the mean-strain time. For \emph{axisymmetric contraction}
they find that the anisotropy evolution depends only on the total strain $c$ and
is essentially independent of $S^\ast$; this makes the axisymmetric case a clean,
parameter-free target, since matching $c$---a purely kinematic quantity that the
dilatation reproduces exactly---is sufficient. For \emph{plane strain} the
situation is more nuanced: the compressed ($b_{22}$) and stretched ($b_{33}$)
components are nearly $S^\ast$-independent, but the unstrained component $b_{11}$
is strongly $S^\ast$-dependent, ranging in the DNS from $b_{11}\approx0$ at the
lowest strain rates (run~PXH) to the rapid-distortion value $b_{11}\approx+0.12$
at the highest (run~PXG). This distinction is central to the interpretation of
the plane-strain comparison below.

\subsubsection{Plane strain}
\label{sec:lr-plane}
\begin{figure}
  \centering
  \includegraphics[width=0.8\textwidth]{Figures/fig_LR_plane.pdf}
  \caption{Reynolds-stress anisotropy under plane strain: present model
  (lines; ensemble mean over $N=1000$ realizations, standard-error band shown)
  versus the DNS of \citet{LeeReynolds1985} (symbols, digitised from their
  Figure~5.4). The compressed component $b_{22}$ and stretched component $b_{33}$
  are reproduced quantitatively. The unstrained component $b_{11}$ is strongly
  strain-rate-dependent in the DNS; its range from slow to fast strain rate is
  shown as the shaded band, with the model's near-zero value lying at the
  slow-strain edge. Reference total strain $c=\exp\int S\,\mathrm{d}t$.}
  \label{fig:lr-plane}
\end{figure}
Figure~\ref{fig:lr-plane} shows the evolution of the three anisotropy components
under plane strain, with the present model (lines, ensemble mean over $N=1000$
realizations with the standard-error band shown) overlaid on the DNS of
\citet{LeeReynolds1985} (symbols, digitised from their Figure~5.4).

The agreement on the two strained components is good. The compressed-direction
anisotropy $b_{22}$ rises monotonically and the model tracks the DNS closely
across the entire strain range, reaching $b_{22}=+0.21$ at $c=3.86$ against the
DNS value of $\approx+0.20$. The stretched-direction anisotropy $b_{33}$ falls to
$b_{33}=-0.20$ in the model against $\approx-0.26$ in the DNS; the model captures
the correct trend and the bulk of the magnitude, under-predicting the depletion of
the stretched component by about $0.05$ at the largest strain. These two
components are precisely the ones \citeauthor{LeeReynolds1985} identify as
nearly strain-rate-independent, so they constitute the robust part of the
comparison, and the model reproduces them quantitatively.

The unstrained-direction anisotropy $b_{11}$ requires more careful
interpretation, and we present it honestly rather than as a clean match. In the
DNS, $b_{11}$ is the most strain-rate-sensitive component, spanning the shaded
band in Figure~\ref{fig:lr-plane} from $b_{11}\approx0$ (slow strain) to
$b_{11}\approx+0.12$ (fast strain, approaching the rapid-distortion limit). The
present model yields $b_{11}\approx-0.01$, essentially zero, whereas the DNS
rises to $b_{11}\approx+0.12$ at the highest strain rate. As shown in
\S\,\ref{sec:lr-discussion}, this is not a deficiency of the present model: the
positive $b_{11}$ seen in the DNS lies outside the reach of rapid-distortion
theory and of the standard second-moment closures alike, and the present result
coincides with the closure prediction on which the model is built. Physically,
$b_{11}$ is governed entirely by intercomponent redistribution, since the
unstrained direction receives no energy directly from the mean strain; we
therefore defer its interpretation to \S\,\ref{sec:lr-discussion}, where it is
compared against the full hierarchy of analytic predictions.

A consistency check supports the internal correctness of the computed
anisotropy: the model's $b_{ij}$ is traceless to machine precision at every
strain, $b_{11}+b_{22}+b_{33}=0$, as required by the definition~\eqref{eq:bij}.

\subsubsection{Axisymmetric contraction}
\label{sec:lr-axisym}
\begin{figure}
  \centering
  \includegraphics[width=0.8\textwidth]{Figures/fig_LR_axisym.pdf}
  \caption{Reynolds-stress anisotropy under axisymmetric contraction: present
  model (lines; $N=1000$) versus the DNS of \citet{LeeReynolds1985} (symbols,
  digitised from their Figure~5.8). The axis component $b_{11}$ and the transverse
  mean $\tfrac{1}{2}(b_{22}+b_{33})$ are compared; the shaded spread is the
  residual $b_{22}\neq b_{33}$ split arising because the one-dimensional line
  breaks the transverse symmetry of the flow. The model reproduces the trend and
  sign structure while under-predicting the anisotropy magnitude by about
  $25\%$.}
  \label{fig:lr-axisym}
\end{figure}
Figure~\ref{fig:lr-axisym} shows the axisymmetric-contraction comparison. Here
the physically distinct quantities are the axis component $b_{11}$ (the stretched
direction) and the transverse anisotropy of the two compressed directions, which
in the exact flow are equal, $b_{22}=b_{33}$, by the axisymmetry of the imposed
strain about $x_1$.

A structural feature of the one-dimensional representation must be stated at the
outset. Because the ODT line lies along one of the two transverse (compressed)
directions, the model resolves the dynamics along that direction explicitly
through the eddy events, whereas the orthogonal transverse direction is
represented only through the pointwise three-component velocity. The line
therefore breaks the $x_2\leftrightarrow x_3$ symmetry that the physical flow
possesses, and the model does not return exactly $b_{22}=b_{33}$. At $N=1000$
realizations the residual split is small---$b_{22}=+0.102$ versus
$b_{33}=+0.113$ at $c=3.16$, a difference of about $0.011$---and is shown as the
shaded spread in Figure~\ref{fig:lr-axisym}; the larger asymmetry seen in a
single realization is sampling noise that the ensemble removes. The
symmetry-respecting transverse quantity is the mean
$\tfrac{1}{2}(b_{22}+b_{33})$, which by tracelessness equals $-\tfrac{1}{2}b_{11}$
and is the appropriate object to compare against the DNS transverse curve.

With this understood, the model reproduces the correct qualitative structure: the
stretched axis loses anisotropy ($b_{11}<0$) while the compressed transverse
plane gains ($b_{22},b_{33}>0$), monotonically with total strain. Quantitatively,
the model under-predicts the magnitude of the anisotropy: at $c=3.16$ it reaches
$b_{11}=-0.21$ against the DNS value of $\approx-0.29$ (digitised from their
Figure~5.8), an under-prediction of roughly $25\%$, with a correspondingly
smaller transverse anisotropy. The trend, the sign structure, and the monotonic
approach to a strained asymptotic state are all captured; the magnitude is
systematically low.

\subsubsection{Discussion: the unstrained component and the limits of the closure
family}
\label{sec:lr-discussion}
\begin{figure}
  \centering
  \includegraphics[width=0.82\textwidth]{Figures/fig_LR_b11_closure.pdf}
  \caption{The unstrained-direction anisotropy $b_{11}$ under plane strain: the
  present ODT (solid, $N=1000$, with standard-error band) compared with exact
  linear rapid-distortion theory (dotted), the isotropization-of-production
  closure (dashed), the quasi-isotropic Launder--Reece--Rodi closure (the form
  coded in the model), and the DNS of \citet{LeeReynolds1985} (symbols and shaded
  $S^\ast$ range). The present model coincides with the LRR-QI closure prediction;
  exact rapid distortion gives $b_{11}<0$, since the unstrained component is
  conserved while the energy grows. The DNS rise to $b_{11}>0$ lies above the
  entire closure family and reflects finite-Reynolds-number vorticity structure
  outside the moment-closure framework.}
  \label{fig:lr-b11}
\end{figure}
Taken together, the two strain modes establish that the model reproduces the
Reynolds-stress anisotropy response of strained homogeneous turbulence at the
moment level---quantitatively for the production- and rapid-distortion-governed
components $b_{22}$ and $b_{33}$, and qualitatively in all cases---against DNS at
a Reynolds number that, unlike the spectral benchmark of \S\,\ref{sec:cmk}, is
fully resolved on the single line.

The one component the model does not reproduce, the plane-strain unstrained
anisotropy $b_{11}$, deserves a careful and honest interpretation, because it is
\emph{not} a deficiency specific to the present model: it lies outside what the
rapid-distortion limit and the entire second-moment-closure family can produce.
This is established directly in Figure~\ref{fig:lr-b11}, which compares the
present ODT against the analytic alternatives on the same axes. Exact linear
rapid-distortion theory gives $b_{11}<0$ in plane strain, and the reason is
kinematic and exact: the unstrained direction has $A_{1k}=0$, so under rapid
distortion $\langle u_1^2\rangle$ is conserved while the total energy $q^2$ grows,
forcing $b_{11}=\langle u_1^2\rangle/q^2-\tfrac{1}{3}$ to \emph{decrease}.
Integrating the standard second-moment closures over the same strain---both the
quasi-isotropic Launder--Reece--Rodi form coded in the model and the
isotropization-of-production form---yields $b_{11}\approx0$, slightly negative,
for the full trajectory. The present ODT result coincides with the LRR-QI closure
curve to within the ensemble scatter. In other words, the model returns exactly
the prediction of the second-moment closure on which its strain operator is built;
its $b_{11}$ is correct \emph{for its formulation}.

The DNS, by contrast, shows $b_{11}$ rising to $+0.12$ at high strain rate. That
the effect is largest at \emph{fast} strain rules out a slow return-to-isotropy
origin, which would require time to act and would therefore be largest at slow
strain; the positive $b_{11}$ is a rapid effect that nonetheless lies beyond the
reach of rapid-distortion theory and of the standard rapid pressure--strain
closures. \citet{LeeReynolds1985} attribute the velocity-field anisotropy in
these flows to the structure of the turbulent vorticity field---an essentially
nonlinear, finite-Reynolds-number property that second-moment closures, which
carry no vorticity-anisotropy information, do not represent. The departure of the
DNS $b_{11}$ from the entire closure family in Figure~\ref{fig:lr-b11} is thus a
known limitation of moment-level modelling, shared by rapid-distortion theory and
by every closure shown, rather than an artifact of the present approach.

The practical consequence for the present work is favourable. The components that
govern the upwash spectrum fed to the leading-edge response---the energy
amplification and the anisotropy in the strained directions---are precisely those
the model reproduces quantitatively. The single component it does not capture is
one that no rapid-distortion or second-moment closure captures, and whose
finite-Reynolds-number origin is outside the modelling framework adopted here. The
validation therefore places the model on a firm footing: it performs at the level
of the best available second-moment closure on this benchmark, with the residual
discrepancy attributable to physics that the closure family as a whole omits.

Finally, we note that \citeauthor{LeeReynolds1985}'s own transport-budget analysis
(their Figures~5.5 and~5.6) provides independent support for a central premise of
\S\,\ref{sec:gate}: in plane strain they find that the slow pressure--strain-rate
term is negligible and that the intercomponent transfer is carried almost entirely
by the rapid term. This is the same decomposition on which the closure of
\S\,\ref{sec:gate} rests, and the DNS confirmation that the rapid term dominates in
this flow lends physical weight to the closure beyond the anisotropy comparison
itself.








\bibliographystyle{jfm}
\bibliography{jfm}


\end{document}